% concatenated from log_lin_conv_arxiv.tex
\documentclass{amsart}
\usepackage[a4paper, margin=1in]{geometry}

\usepackage[utf8]{inputenc} 
\usepackage[T1]{fontenc}    
\usepackage{hyperref}       
\usepackage{url}            
\usepackage{booktabs}       
\usepackage{amsfonts}       
\usepackage{nicefrac}       
\usepackage{microtype}      
\usepackage{xcolor}         

\usepackage{graphicx}       
\usepackage{subcaption}      
\usepackage{thmtools}        
\usepackage{thm-restate}    

\usepackage[round, authoryear]{natbib}


\newtheorem{theorem}{Theorem}                   
\newtheorem{lemma}[theorem]{Lemma}              
\newtheorem{conjecture}[theorem]{Conjecture}    
\newtheorem{definition}{Definition}             
\newtheorem{proposition}{Proposition}           
\newtheorem{assumption}{Assumption}             
\newtheorem{corollary}{Corollary}               

\theoremstyle{definition} 
\newtheorem{example}{Example}                 
\newtheorem{remark}{Remark}                     

\title[Global linear convergence of softmax policy gradient]{Global linear convergence of entropy-regularized softmax policy gradient beyond tabular MDPs}


\author{Ziyue Chen}
\address{School of Mathematics\\University of Edinburgh\\Edinburgh, UK, EH9 3FD}
\email{s2674476@ed.ac.uk}

\author{David {\v{S}}i{\v{s}}ka}
\address{School of Mathematics\\University of Edinburgh\\Edinburgh, UK, EH9 3FD}
\email{d.siska@ed.ac.uk}

\author{{\L}ukasz Szpruch}
\address{School of Mathematics\\University of Edinburgh\\Edinburgh, UK, EH9 3FD}
\email{L.Szpruch@ed.ac.uk}

\begin{document}


\begin{abstract}
We study the global convergence of policy gradient for infinite-horizon entropy-regularized Markov decision processes (MDPs) with continuous state and action spaces.
We consider log-linear softmax policies with linear function approximation, which extend the tabular softmax parameterization while retaining a tractable  policy class. Under $Q^\pi_\tau$-realizability for the regularized state-action value function, we first establish a non-uniform Polyak--{\L}ojasiewicz (P\L) inequality.
The non-uniformity arises through degeneracy of constants associated with the policy geometry, namely the Fisher information matrix or an uncentered feature covariance matrix. We then identify two feature regimes under which this non-uniform constant can be bounded along the gradient flow.
For full-affine-span features, we prove radial unboundedness of the KL regularizer and show that the smallest eigenvalue of the Fisher information matrix remains bounded below by an initialization-dependent positive constant.
For simplex-valued features, we prove an analogous radial unboundedness result in the subspace orthogonal to the all-ones vector and obtain a uniform lower bound for the smallest eigenvalue of the uncentered covariance matrix.
These results imply global linear convergence of the regularized objective along the gradient flow, i.e. suboptimality decaying as $\mathcal{O}(e^{-Ct})$ for some $C>0$.
Our analysis extends the global convergence theory of entropy-regularized softmax policy gradient beyond the tabular setting of~\cite{agarwal2020optimality, bhandari2024global, mei2020global}.
\end{abstract}

\maketitle


\section{Introduction}
\label{sec:intro}

\textbf{Overview:} The policy gradient is a fundamental concept in reinforcement learning (RL), underpinning policy search and actor-critic methods \cite{sutton1999policy}. However, for parametrized softmax policies, the normalization factor induces non-convexity in the parameters. As a result, despite the prevalence of softmax policies in RL, their theoretical understanding has remained limited until recently, with existing results confined to the tabular setting \cite{agarwal2020optimality,agarwal2021theory,bhandari2024global,mei2020global}.

We consider a discounted infinite horizon Markov decision model $(S,A,P,c,\gamma)$, where $S$ and $A$ are general, possibly continuous, state space and action spaces respectively.
Let $\mu$ be the fixed finite reference measure, $\rho$ the distribution of initial state, $P\in \mathcal{P}(S|S\times A)$ the transition probability kernel, $c$ a bounded cost function, and $\gamma\in [0,1)$ the discount factor.
For a given stochastic policy $\pi \in \mathcal{P}(A|S)$, we define the entropy regularized value function $V^{\pi}_{\tau}:S\rightarrow \mathbb{R}$  by
\begin{equation}\label{eq:intro:objective}
V^{\pi}_{\tau}(\rho)= \mathbb{E}_{s_0 \sim \rho,s_{n+1}\sim P(\cdot|s_n,a_n)}^{a_n\sim \pi(\cdot|s_n)}\left[\sum_{n=0}^{\infty}\gamma^{n} \left(c(s_n,a_n)+\tau\,\log\tfrac{d\pi}{d\mu}(a_n|s_n)\right)\right],
\end{equation}
where $\tau > 0$ determines the intensity of the entropy regularization.
For full details on our assumptions and notation, we refer to Section \ref{sec:nota and def}.

There are two implications of having $\tau>0$.
The first is that the optimal policy satisfies
\begin{equation}\label
{eq:intro:optimal_policy}
\pi_\tau^*(d a|s)=\exp \left(\tfrac{1}{\tau}\left(Q_\tau^*(s, a)-V_\tau^*(s)\right)\right) \mu(d a),
\end{equation}
where $V_\tau^*$ and $Q_\tau^*$ denote the (bounded) optimal value and state-action value functions, respectively. This is an immediate consequence of the Bellman principle (see Theorem \ref{thm:app:dpp} or ~\cite{ziebart2010modeling, haarnoja2017reinforcement, geist2019theory}). 

Second, since the entropy term is strictly convex~\cite[Section 1.4]{dupuis2011weak}, its addition in~\eqref{eq:intro:objective} is expected to improve the convergence when optimizing $V_\tau^\pi(\rho)$ over $\pi$ using policy gradient.
While the latter point may seem intuitive, the analysis is far from being straightforward even in the tabular case when $S$ and $A$ are finite and
direct parametrization 
\begin{equation}\label{eq:softmax}  
\pi_\theta(a|s)\propto \exp \theta(s,a)
\end{equation}
with $\theta: S\times A\rightarrow \mathbb{R}$ is employed.
Two main difficulties arise: first $\theta \in \mathbb{R}^p \mapsto V^{\pi_\theta}_\tau(\rho)$ is non-convex (see, e.g., Proposition 1 in~\cite{mei2020global}), even in the bandit case. 
Moreover, second, $\mathcal P(A|S) \ni \pi \mapsto V_\tau^\pi(\rho)$ is in general non-convex~\cite{agarwal2021theory, Giegrich2024} even when dynamics are linear and costs convex.
Nevertheless, convergence with good rates of policy gradient with softmax policies in the tabular setting has been shown in~\cite{mei2020global}.
In general, the suboptimality $V^{\pi_{\theta_t}}_0 - V^\ast_0$ converges sub-linearly, i.e. $\mathcal O(1/t)$, while with the additional entropy regularization, the the suboptimality $V^{\pi_{\theta_t}}_\tau - V^\ast_\tau$ converges converges linearly, i.e. $\mathcal O(e^{-Ct})$.
The key insight in~\cite{mei2020global}, is to use non-uniform Polyak--\L ojasiewicz (P\L) inequality.

This approach becomes computationally intractable as the size of the sets $S$ and $A$ grow large or when $S$ or $A$ are continous. 
To overcome this one parametrizes the log densities. 
In this paper we assume the following linear function approximation: given basis functions $g: S \times A \rightarrow$ $\mathbb{R}^p$, for all $a \in A$ and $s \in S$, let
\begin{equation*}
\pi_\theta(da|s) \propto \exp \left(\langle\theta,g(s,a)\rangle\right)\mu(da),
\end{equation*}
where $\theta \in \mathbb{R}^p$. 
The continuous-time policy gradient is 
\begin{equation}\label
{eq:intro:continuous_GF}
\frac{d}{dt}\theta_t = - \nabla_{\theta} V_\tau^{\pi_{\theta_t}}(\rho)\,,\,\,\, t \geq 0\,,\,\,\, \theta_0 \in \mathbb R^p\,\,\,\text{given}.
\end{equation}
The expression for the gradient is given by the well known policy gradient theorem, which we restate for convenience later, in Proposition~\ref{prp:app:mdp:Gradient_of_the_Objective}.


\textbf{Outline of the argument to obtain linear convergence:} 
First, we wish to obtain gradient dominance property for the objective in the form of a P{\L} inequality.
To proceed we assume $Q^\pi_\tau$-realizability i.e. that for any $\pi$ there exists a unique $\boldsymbol{{\theta}}(\pi)$ such that $\langle \boldsymbol{{\theta}}(\pi), g \rangle = -\tfrac{1}{\tau} Q^{\pi}_\tau$. 
Under this assumption a simple calculation shows that 
\begin{equation}
\label{eq:pl_calc0_intro}
\nabla_\theta V^{\pi_\theta}_\tau(\rho) = -\tfrac{\tau}{1-\gamma} \Big( \textstyle\int_S G^{\pi_\theta}(s)d^{\pi_\theta}_\rho(ds)\Big)(\theta - \boldsymbol{{\theta}}(\pi_\theta))\,,
\end{equation}
where the Fisher Information matrix (FIM) of $\pi_\theta$ for fixed $s$ is defined as
\begin{equation}\label
{eq:intro:FIM}
G^{\pi_\theta}(s) = \textstyle\int_A \nabla_\theta \log \tfrac{d \pi_{\theta}}{d \mu}(a|s)\big(\nabla_\theta \log \tfrac{d \pi_{\theta}}{d \mu}(a|s)\big)^{\top}\pi_{\theta}(da|s)\,.
\end{equation}
Left multiplying~\eqref{eq:pl_calc0_intro} by $(\theta-\boldsymbol{\theta}(\pi_\theta))^\top$, using that the FIM is positive semi-definite, using the Cauchy--Schwartz inequalty and finally dividing by $\|\theta-\boldsymbol{\theta}(\pi_\theta)\|_2$ we have
\begin{equation}
\label{eq:pl_calc1_intro}
\|\nabla_\theta V^{\pi_\theta}_\tau(\rho)\|_2 \geq \tfrac{\tau}{1-\gamma}  \textstyle\int_S \lambda_{\text{min}}(G^{\pi_\theta}(s)) d^{\pi_\theta}_\rho(ds)\|\theta - \boldsymbol{{\theta}}(\pi)\|_2\,.
\end{equation}
Separately, using the KL-sandwich inequality (see Lemma~\ref{lem:sandwich-inequality}) and local Lipschitz continuity of KL (see Lemma~\ref{lem:app:kl-logit_ineq_affine}, which is an extension of the tabular case from~\cite[Lemma 27]{mei2020global}) we get that for some $C_{\tau,\gamma,\theta}>0$
\[
0 \leq V_\tau^{\pi_{\theta}}(\rho)-V_\tau^{\pi^*}(\rho) \leq C_{\tau,\gamma,\theta} \|\theta - \boldsymbol{{\theta}}(\pi)\|_2^2\,.
\]
This, together with~\eqref{eq:pl_calc1_intro} 
show that there is $ C: \mathbb{R}^p \to (0, \infty)$ s.t. for all
$\theta \in \mathbb{R}^p$
\begin{equation}
\label{eq:pl_ineq_intro}
0 \leq V_\tau^{\pi_{\theta}}(\rho)-V_\tau^{\pi^*}(\rho) \leq  C(\theta)\|\nabla_\theta V^{\pi_\theta}_\tau(\rho)\|_2^2 \,.
\end{equation}
This is stated fully as Theorem~\ref{thm:main: non-uniform_PL_ineq} below with a complete proof given later.
Note that the strength of the gradient dominance is parameter dependent i.e. we only have non-uniform P{\L} inequality.


Neverthless with $\{\theta_t\}_{t\geq 0}$ given by the gradient flow~\eqref{eq:intro:continuous_GF} we have
\[
\tfrac{d}{d t}\Big[V_\tau^{\pi_{\theta_t}}(\rho)-V_\tau^{\pi^*}(\rho)\Big]=-\left\|\nabla V_\tau^{\pi_{\theta_t}}(\rho))\right\|^2 \leq- C^{-1}\left(\theta_t\right)\Big[V_\tau^{\pi_{\theta_t}}(\rho)-V_\tau^{\pi^*}(\rho)\Big]\,.
\]
Hence, from Gr\"onwall's lemma we immediately get that
\[\
0 \leq V_\tau^{\pi_{\theta_t}}(\rho)-V_\tau^{\pi^*}(\rho) \leq\big[V_\tau^{\pi_{\theta_0}}(\rho)-V_\tau^{\pi^*}(\rho)\big] \exp \left(-\textstyle\int_0^t  C^{-1}\big(\theta_s\right) d t\big) \,.
\]
If we can show that non-uniform term can be lower bounded along the flow, i.e. if we can show that $\inf _t C^{-1}\left(\theta_t\right) > 0$ then this is the required linear convergence rate.

We notice that if $\{\theta_t\}_{t\geq 0}$ given by the gradient flow~\eqref{eq:intro:continuous_GF} produces policies with uniformly bounded log densities i.e. if $\sup_{t} |\log \frac{\mathrm d \pi_{\theta_t}}{\mathrm d\mu}|_{B_b(S\times A)} < \infty$ then the smallest eigenvalue of the FIM will remain bounded away from zero and so we will have $\inf_t C^{-1}\left(\theta_t\right) > 0$.
Thus it is enough to show that $\{\theta_t\}_{t\geq 0}$ is contained in a compact subset of $\mathbb R^p$.
We will use Lyapunov function techniques to that end.

With that in mind, we note that a simple calculation using the chain rule yields value improvement along the gradient flow~\eqref{eq:intro:continuous_GF}, namely
$
\frac{d}{dt}V^{\pi_{\theta_t}}_\tau(\rho) = - \|\nabla_\theta V^{\pi_\theta}_\tau(\rho)\|_2^2 \leq 0\,.
$
The value function $\theta \mapsto V^{\pi_{\theta_t}}_\tau(\rho)$ can then be used as a Lyapunov function as long as it is radially unbounded.
Since the cost itself is bounded the radial unboundedness can only come as a consequence of the KL term. 
Example~\ref{ex:hat_func_short}, which uses the hat functions as a basis (P1 finite elements) shows that not every reasonable feature basis leads to a radially unbounded KL term.
Howevever we identify conditions on the feature basis which ensure radial unboundedness i.e. that $\operatorname{KL}(\pi_\theta|\mu) \to \infty$ whenever $\|\theta\|_2 \to \infty$.
In Example~\ref{ex:trig_func_mdp} we show that the Fourier basis satisfies our Assumption~\ref{asp:app:mdp:max_set_affine}.

Additionally, in the case of simplex features (see Assumption~\ref{asp:app:mdp:max_set_simplex}) below we can extend the P{\L} inequality by replacing the FIM with a version built using uncentred features which allows us to carry out a similar convergence argument but capturing further types of basis functions, e.g. the Bernstein polynomials, see Example~\ref{ex:Bernstein_Poly_MDP}.

\textbf{Key differences with the tabular case:}
In the tabular case when $\tau > 0$~\cite{mei2020global} prove P{\L} inequality~\eqref{eq:pl_ineq_intro} with $C(\theta) = C \min_{s,a} \pi_{\theta_t}(a|s)$ with $C$ independent of $\theta$.  
Their methods then can be used to show that along the flow~\eqref{eq:intro:continuous_GF} it holds that 
$\sum_a \tfrac{\partial V^{\pi_{\theta_t}}_\tau(\rho)}{\partial \theta_t(s, a)} =0$ for all $t$. 
From this they then derive $\inf_t \min_a \pi_{\theta_t}(a) > 0$ which, as we've seen, leads to linear convergence.

In our setting, with general log-linear policies, their way to formulate the property that the gradient of the value function summed up over $a$ is $0$ along the flow cannot be employed. 
Even though with simplex features we have $\sum_{i} \nabla_{\theta_i} V^{\pi_{\theta_t}}_\tau(\rho) = 0$ for all $t$, weights $\theta$ impact the entire conditional log-density and thus we cannot hope to derive a lower bound for the  action density at a given state since this separation is not available.
Thus, a fundamentally different approach has to be used.
On the other hand, in the tabular setting, one cannot expect radial unboundedness of the KL divergence as this term will be finite any $\mu$ s.t. $\mu(a) > 0$ for all $a$. 
Thus the use of the KL divergence as a Lyapunov function is novel and our results complement those of~\cite{mei2020global}.


\textbf{Literature review:} 
There is a tremendous amount of research literature on convergence RL methods, which underscores its importance. 
Here we focus on the subset of the RL literature that we think is most related to our work.

Entropy-regularized RL has demonstrated both good algorithmic performance and desirable theoretical properties~\cite{haarnoja2017reinforcement,haarnoja2018soft, geist2019theory,vieillard2020leverage, neu2017unified, fox2015taming,ziebart2010modeling}. 
It has been shown that the softmax policies are optimal in the entropy regularized setting.

For policy gradient in tabular setting, the work of \cite{agarwal2020optimality} initiated a global analysis in discounted MDPs under tabular setting and compatible function approximation parameterizations, making explicit the roles of distribution mismatch, approximation error, and statistical error; in the tabular softmax case this gives a sublinear convergence guarantee for vanilla policy gradient. 
The work of \cite{mei2020global} subsequently gave a sharper analysis of tabular softmax policy gradient with direct parametrization as discussed above. 
These guarantees are complemented by worse case analysis for the policy gradient methods in~\cite{agarwal2020optimality, mei2020global} and show that standard softmax policy gradient the constants in the convergence rate suffer from exponential dependence on state-space and horizon sizes~\cite{li2021softmax}.
Moving beyond the realizability assumption~\cite{lin2025rethinking} derive some  ordering conditions in the bandit case (empty state space) which guarantee the convergence of softmax policy gradient using linear function approximation.
However as of now it is unclear how this applies to MDPs and also fundamentally depends on the finite cardinality of the action space.

Although MDPs with continuous state and action spaces are widely used in practical applications \cite{doya2000reinforcement, van2012reinforcement, manna2022learning}, the convergence analysis of policy gradient methods in this setting remains less developed compared to its discrete counterparts.
Most existing works have focused on discrete-time linear quadratic regulator (LQR) problems with linear parameterized policies.
The linear-quadratic structure leads to the P{\L} inequality~\cite{polyak1963gradient, lojasiewicz1963topological, kurdyka1998gradients} with a uniform constant and we've seen above how this then gives linear convergence~\cite{fazel2018global, bu2019lqr, hu2023toward}.
These results have been extended to continuous-time LQR systems~\cite{sontag2022remarks,Giegrich2024}. 

One may choose to use the inverse of the integrated FIM~\eqref{eq:intro:FIM} as a pre-conditioner in the gradient flow~\eqref{eq:intro:continuous_GF} which gives rise to the natural policy gradient (NPG).  
For NPG under with log-linear policies~\cite{cayci2024convergence} prove linear convergence of the NPG in the entropy regularized setting.
Moreover for log-linear policies the NPG is in fact identical to mirror descent (MD) (in the sense of producing the same policy updates) and MD is known to converge linearly for entropy regularized MDPs~\cite{lan2023policy, ju2022policy, kerimkulov2025fisher}.
It is perhaps interesting to note that unlike policy gradient, MD / NPG automatically ensure that log densities remain bounded along the optimization and thus the FIM is invertible almost regardless of the feature basis (one cannot have linearly dependent features).
This is a consequnce of the policy improvement of NPG / MD under exact evaluations. 

A linear convergence rate is proved in \cite{liu2023polyak} under an additional P{\L} inequality for the continuous-time Fisher--Rao flow on the space of measures.
Continuous-time Fisher--Rao flows in the entropy regularized MDPs has been studied by~\cite{kerimkulov2025fisher}, where the linear convergence to the optimal policy has been established and the insights into the natural policy gradient flow with linear function approximation was given. 

\textbf{Main contributions of this paper:}
We prove that under suitable conditions on the features and under $Q^\pi_\tau$-realizability the policy gradient for the KL-regularized MDP on general state and action spaces with exact evaluations converges linearly. 
In particular we:
\begin{enumerate}
\item Prove a non-uniform P{\L} inequality in this setting, see Theorem~\ref{thm:main: non-uniform_PL_ineq}.
\item Establish conditions on the feature basis under which the KL divergence is radially unbounded and hence Lyapunov function techniques can be used to obtain a bound on the non-uniform constant, see Theorems~\ref{thm:main:mdp:Radial_unboundedness_of_relative_entropy} and~\ref{thm:main:mdp:Radial_unboundedness_of_relative_entropy_in_subspace_simplex}.
This is of independent interest to anyone analysing algorithms that employ log-linear densities and KL regularization as it opens up Lyapunov function techniques to them.
\item Obtain the desired linear convergence, see Theorems~\ref{thm:main:lin_cov_affine} and~\ref{thm:main:lin_cov_simplex}.
\item Provide examples of basis functions which satisfy the conditions, see Examples~\ref{ex:trig_func_mdp} and~\ref{ex:Bernstein_Poly_MDP}. 
\end{enumerate}

\section{Formulation and the statement of the main results}

Let $S$ and $A$ be Polish spaces with $A$ compact.
Let $P\in \mathcal{P}(S|S\times A)$.
Let  $\gamma\in [0,1)$ be the discount factor.
Let $c\in B_b(S\times A)$ be the cost function and we will assume, without loss of generality, that  $|c|_{B_b(S \times A)}\leq 1$. 
Let $\tau > 0$. 
Let $\mu \in \mathcal P(A)$ have full support on $A$ and be absolutely continuous w.r.t. the Lebesgue measure.
The seven-tuple $(S,A,P,c,\gamma,\tau, \mu)$ determines a $\gamma$-discounted infinite horizon $\tau$-entropy regularized Markov decision process model.

For a given randomized Markov policy $\pi\in P(A|S)$, we define the $\tau$-entropy regularized value function $V^{\pi}_{\tau}:S\rightarrow \mathbb{R}\cup\{+\infty\}$  by
\[
V^{\pi}_{\tau}(\rho)= \mathbb{E}_{s_0 \sim \rho,s_{n+1}\sim P(\cdot|s_n,a_n)}^{a_n\sim \pi(\cdot|s_n)}\left[\sum_{n=0}^{\infty}\gamma^{n} \left(c(s_n,a_n)+\tau\,\log\tfrac{d\pi}{d\mu}(a_n|s_n)\right)\right].
\]
The aim is to minimize $P(A|S)\ni \pi \mapsto V^{\pi}_{\tau}(\rho)$.
Due to the Bellman principle, Theorem \ref{thm:app:dpp} for convenience in Section \ref{sec:extra results} or ~\cite{ziebart2010modeling, haarnoja2017reinforcement, geist2019theory}, 
we know that the optimal policy is of the form~\eqref{eq:intro:optimal_policy} 
and the optimal value (and state value) functions are bounded.
Hence, without loss of generality we may restrict our minimization problem to softmax policies from the class
\begin{align}\label
{def:Pi_class}
\Pi_{\mu} &=\left\{\tfrac{\exp(f(s,a))\mu(da)}{\int_{A} \exp(f(s,a^\prime))\mu(da^\prime)}\in \mathcal{P}_{\mu}(A|S) | f\in B_b(S\times A)\right\}  \,.
\end{align}



For a given policy $\pi\in\Pi_{\mu}$, we define the regularized state-action value function $Q^{\pi}_{\tau}\in B_b(S\times A)$ by
\begin{equation}\label{def:intro:Qpi}
Q^{\pi}_{\tau}(s,a)=c(s,a)+\gamma\int_{S}V^{\pi}_{\tau}(s')P(ds'|s,a)\,.
\end{equation}

The occupancy kernel $d^{\pi}\in \mathcal{P}(S|S)$ is defined by $d^{\pi}(ds'|s)=(1-\gamma)\sum_{t=0}^{\infty}\gamma^t P^t_{\pi}(ds'|s)\,, d^{\pi}_{\rho}(ds)=\int_{S}d^{\pi}(ds|s')\rho(ds')$
where $P^0_{\pi}(ds'|s):=\delta_s(ds')$, $P^t_{\pi}$ is understood as a product of kernels, and convergence is understood in $b\mathcal{K}(S|S)$. 
It is well known that the on policy Bellman equation holds, see e.g.~\cite[Lemma B.2] {kerimkulov2025fisher}, that is
\begin{equation}
\label{eq:intro:policy_Bellman}
V^{\pi}_{\tau}(s) = \int_{A}\left(c(s,a)+\tau\log\frac{d\pi}{d\mu}(a|s) +\gamma\int_{S}V^{\pi}_{\tau}(s')P(ds'|s,a)\right)\pi(da|s) \,\,\, \forall s\in S\,.
\end{equation}
Moreover one can see that this has the stochastic representation 
\begin{equation}
\label{eq:intro:value_function_occ}
V^{\pi}_{\tau}(s) =\frac{1}{1-\gamma}\int_{S}\int_{A}\left(c(s',a')+\tau \log \frac{d\pi}{d\mu}(a'|s')\right)\pi(da'|s')d^{\pi}(ds'|s)\,\,\, \forall s\in S\,.
\end{equation}
For a given fixed initial distribution $\rho\in\mathcal{P}(S)$, we define 
\[
V^{\pi}_{\tau}(\rho)=\int_{S}V^{\pi}_{\tau}(s)\rho(ds), \;  d^{\pi}_{\rho}(ds)=\int_{S}d^{\pi}(ds|s')\rho(ds')\,.
\]

When $S$ or $A$ are not of finite cardinality the minimization over the class~\ref{def:Pi_class} is intractable. 
Thus, instead of looking for the optimal policy in~\ref{def:Pi_class}, we parametrize the softmax policies using linear function approximation.
To that end let $g: A \times S \rightarrow$ $\mathbb{R}^p$ be our basis functions (or features). 
We will take $g$ to be measurable and such that $\sum_{i=1}^p|g_i|_{B_b(S\times A)}^2 \leq 1$.
Now given parameters $\theta \in \mathbb R^p$ consider parametrized policies of the form
\[
\pi_\theta(da|\cdot)=\tfrac{\exp \left(\langle\theta,g(\cdot,a)\rangle\right)\mu(da)}{\int_{A} \exp \left(\langle\theta,g(\cdot,a^\prime)\rangle\right)\mu(da^\prime)}.
\]
Thus we wish to solve the minimization problem
\begin{equation}\label{eq:intro:entropy regularized mdp0}
\min_{\theta \in \mathbb{R}^p} V_\tau^{\pi_\theta}(\rho)\,.
\end{equation}
We will work under the $Q^\pi_\tau$-realizability assumption.

\begin{restatable}[$Q^\pi_\tau$-realizability]{assumption}{aspQpiRealizability}\label{asp:app:mdp:Q_realisability}
For any $\pi \in \Pi_\mu$, there exists a unique $\boldsymbol{\theta}(\pi)$ such that $\langle \boldsymbol{\theta}(\pi), g(s,a) \rangle = -\frac{1}{\tau} Q^{\pi}_\tau(s,a)$ for all $s \in S, a \in A$.
\end{restatable}

An example of when this holds besides the tabular case is linear MDPs (see, e.g.,~\cite{yang2019sample, zanette2020frequentist, li2021sample}). 
An MDP is linear if there exists exists $w\in \mathbb{R}^p$ and a sequence $\{\psi_i\}_{i=1}^p$ with $\psi_i \in \mathcal{M}(S)$ such that for all $(s,a)\in S \times A$, $c(s,a)=\langle w,g(s,a)\rangle, P( d s'\mid s,a)=\sum_{i=1}^p g_i(s,a)\psi_i( d s')$.
In this case, given $\pi \in \mathcal P(A|S)$ we can take $\mathbb{\theta}(\pi)_i = -\tfrac1\tau\big(w_i + \gamma\int_{S} V^{\pi}(s')\psi_i(ds')\big)$ 
so that 
\[
\langle \boldsymbol{\theta}(\pi), g(s,a) \rangle = -\tfrac1\tau\langle w, g(s,a)\rangle -\tfrac\gamma\tau \textstyle \int_S V^{\pi}_\tau(s') \sum_{i=1}^p g_i(s,a)\psi_i(ds') = -\tfrac1\tau Q^\pi_\tau(s,a)\,.
\]

Recall that due to the Bellman principle the optimal state-action value function $Q^\ast_\tau = Q^{\pi_\tau^\ast}_\tau \in B_b(S\times A)$ with $\pi^\ast \in \Pi_\mu$.
Let $\theta^*$ such that $\langle \theta^*,g\rangle = -\frac{1}{\tau}Q^{*}_\tau$ which exists and is unique due to Assumption~\ref{asp:app:mdp:Q_realisability}.
Since the optimal $\pi^* \in \Pi_\mu$ is of the form~\eqref{eq:intro:optimal_policy}, we have  
\[\pi^*(da 
| \cdot) \propto \exp\big(-\tfrac{1}{\tau}Q^{\pi^*}_\tau(\cdot,a)\big) \mu(da) = \exp\big(\langle \theta^*, g(\cdot,a) \rangle \big) \mu(da)\,.
\]
In other words, the $Q^\pi_\tau$-realizability assumption for the KL regularized MDP means that our minimization problem is solvable: $\min_{\theta \in \mathbb{R}^p} V_\tau^{\pi_\theta}(\rho) = V^*_\tau(\rho)$.

The remainder of the paper is devoted to the argument that under suitable assumptions on the features the gradient flow~\eqref{eq:intro:continuous_GF} converges linearly in the sense that $0 \leq V_\tau^{\pi_\theta}(\rho) - V_\tau^{\pi_\theta^\ast}(\rho) \leq \mathcal O(e^{-Ct})$.
Recall that we will do this by first obtaining a non-uniform P{\L} inequality and then demonstrating that along the gradient flow we in fact can bound the constant uniformly, using radial unboundedness of $\theta \mapsto \operatorname{KL}(\pi_\theta|\mu)$ which holds for suitable feature basis.

\begin{restatable}[Non-uniform P{\L} \ inequality]{theorem}{PL}\label
{thm:main: non-uniform_PL_ineq}
Let Assumption~\ref{asp:app:mdp:Q_realisability} hold.
Let $R \geq |\log \tfrac{d \pi_{\theta^\ast}}{d\mu}|_{B_b(S\times A)}$.
Let
\[
\Theta_R = \{\theta: \pi_\theta \in \Pi_\mu \ \text{and}\  \vert \log \tfrac{d \pi_\theta}{d\mu}\vert_{B_b(S\times A)} \leq R \}\,.
\]
Then for any $\theta \in \Theta_R$ there exists $C_R(\theta)>0$ such that
\begin{equation}
0\leq V_\tau^{\pi_\theta}(\rho) - V_\tau^{\pi_{\theta^*}}(\rho) \leq C_R(\theta)\left\|\nabla V_\tau^{\pi_\theta}(\rho)\right\|^2_2.
\end{equation}
\end{restatable}
Depending on the assumptions on the feature basis, this will be proved in Section \ref{sec:pf_non_uniform_PL},
where the exact form of the constant $C_R(\theta)$ will be given.
We now need to work with specific assumptions on the basis functions.
Below, we will define ``full affine span features'' and ``simplex features'' and discuss how they allow us to prove the linear convergence.

\subsubsection*{Full affine span features} 
We will say that the feature basis $g:S\times A\to \mathbb R^p$ has full affine dimension if $\operatorname{span}\{g(s,a)-g(s,a'):a,a'\in A \}=\mathbb R^p$.
Recall that $u\in\mathbb S^{p-1}$ if $u\in\mathbb R^p$ and $\|u\|_2 = 1$.

\begin{restatable}[]{assumption}{aspAffineMaxSet}\label{asp:app:mdp:max_set_affine} 
For each fixed
$s\in S$, assume that $a\mapsto g(s,a)\in\mathbb R^p$ is continuous
and that, for every $u\in\mathbb S^{p-1}$,
\[
\mu \Big(\arg\max_{a\in A} u^\top g(s,a)\Big)=0 .
\]
\end{restatable}

We see that Assumption~\ref{asp:app:mdp:max_set_affine} implies 1. in 
Lemma~\ref{rem:app:mdp:affine_span}.
Thus $g(s, A)$ has full affine dimension.
We will work under Assumption~\ref{asp:app:mdp:max_set_affine} as it will be needed later to get the radial unboundedness of $\theta \mapsto \operatorname{KL}(\pi_\theta|\mu)$.



\begin{restatable}[Radial unboundedness of KL divergence]{theorem}{KLUBAFFINE}\label{thm:main:mdp:Radial_unboundedness_of_relative_entropy}
Let Assumption \ref{asp:app:mdp:max_set_affine} hold.
Fix $s\in S$. 
Then
\[
\lim_{\|\theta\|_2\to\infty}
\mathrm{KL}\bigl(\pi_\theta(\cdot\mid s)\,|\,\mu\bigr)=\infty .
\]
\end{restatable}

The proof of Theorem \ref{thm:main:mdp:Radial_unboundedness_of_relative_entropy} can be found in Section \ref{sec:pf_Radial_unboundedness_of_relative_entropy}.

Theorem~\ref{thm:main:mdp:Radial_unboundedness_of_relative_entropy}
combined with classical arguments for constructing ODE solutions with Lyapunov functions will give existence of the solution to the gradient flow~\eqref{eq:intro:continuous_GF}.

\begin{restatable}[Existence of solution to the gradient flow using full affine span features]{lemma}{ExistSolAffine}\label{lem:main:mdp:exist_sol_gf_affine}
Let Assumption~\ref{asp:app:mdp:max_set_affine} hold.
Then there exists solution $\{\theta_t\}_{t \geq 0}$ to~\eqref{eq:intro:continuous_GF}. 
Moreover $\sup_{t \geq 0} \big|\log \tfrac{d \pi_{\theta_t}}{d\mu}\big|_{B_b(S \times A)} < \infty$.
\end{restatable}

The proof of Lemma \ref{lem:main:mdp:exist_sol_gf_affine} can be found in Section \ref{sec:pf_exist_sol_gf_affine}.

\begin{restatable}{assumption}{aspInitAffine}\label{asp:app:mdp:initial_singularity_mdp}
For all $s \in S$, the FIM $G^{\pi_{\theta_0}}(s)$ given by~\eqref{eq:intro:FIM} is positive definite.
\end{restatable}

\begin{restatable}[Linear convergence with full affine span features]{theorem}{LinCovAffine}\label{thm:main:lin_cov_affine}
Let Assumption~\ref{asp:app:mdp:Q_realisability}, \ref{asp:app:mdp:max_set_affine} and \ref{asp:app:mdp:initial_singularity_mdp} hold.
Let $\{\theta_t\}_{t \geq 0}$ be the solution to gradient flow~\eqref{eq:intro:continuous_GF}.
Then there exists $C_{\theta_0} > 0$ such that 
for $C_R(\theta_t)$ from Theorem~\ref{thm:main: non-uniform_PL_ineq} we have $\sup_{t \geq 0} C_R(\theta_t) \leq C_{\theta_0}$ and
\[
0 \leq V_\tau^{\pi_{\theta_t}}(\rho) - V_\tau^{\pi^*}(\rho) \leq e^{-t C_{\theta_0}^{-1}} 
(V_\tau^{\pi_{\theta_0}}(\rho) - V_\tau^{\pi^*}(\rho)).
\]
\end{restatable}
The the exact form of $C_{\theta_0}$ can be found in Section \ref{sec:pf_lin_cov_affine}.

To conclude this part, we introduce an example of full affine span features that satisfy Assumption \ref{asp:app:mdp:max_set_affine}. 

\begin{example}[Trigonometric features]\label{ex:trig_func_mdp}
Let $S\subset\mathbb R^{d_1}$ be a compact state set and let $A\subset\mathbb R^{d_2}$ be a compact action set. 
Let $\mu$ denote the $d_2$-dimensional Lebesgue measure on the action space. 
Choose a finite set $\mathcal K\subset\mathbb Z^{d_2}\setminus\{\mathbf 0\}$ of nonzero action frequencies, with no redundant action modes, for example with only one representative from each pair $\{k,-k\}$. 
For each $k\in\mathcal K$, choose a state-frequency vector $\ell_k\in\mathbb R^{d_1}$ and define the trigonometric feature $g(s,a)$ by using the coordinates $\cos(k^\top a+\ell_k^\top s)$ and $\sin(k^\top a+\ell_k^\top s)$ for $k\in\mathcal K$. 
Thus, for $\theta=(\alpha_k,\beta_k)_{k\in\mathcal K}$, define
$f_\theta(s,a)=\theta^\top g(s,a)=\sum_{k\in\mathcal K}\{\alpha_k\cos(k^\top a+\ell_k^\top s)+\beta_k\sin(k^\top a+\ell_k^\top s)\}$.

The features have no action-constant Fourier mode. 
We assume that the action-frequency has been chosen without redundant modes, so that for every fixed $s\in S$, $\theta\neq0$ implies that the map $a\mapsto f_\theta(s,a)$ is not identically zero. 

For fixed $s$, one has $f_\theta(s,a)=\sum_{k\in\mathcal K}\{\widetilde\alpha_k(s)\cos(k^\top a)+\widetilde\beta_k(s)\sin(k^\top a)\}$, where $\widetilde\alpha_k(s)=\alpha_k\cos(\ell_k^\top s)+\beta_k\sin(\ell_k^\top s)$ and $\widetilde\beta_k(s)=-\alpha_k\sin(\ell_k^\top s)+\beta_k\cos(\ell_k^\top s)$. 
Moreover, $\widetilde\alpha_k(s)^2+\widetilde\beta_k(s)^2=\alpha_k^2+\beta_k^2$. 
Hence, if $\theta\neq \mathbf 0$, then for every fixed $s\in S$ at least one pair $(\widetilde\alpha_k(s),\widetilde\beta_k(s))$ is nonzero.
Since all $k\in\mathcal K$ are nonzero action frequencies and there is no redundant action mode, the fixed-state function $a\mapsto f_\theta(s,a)$ is a non-constant real-analytic function of $a$.

For fixed $s\in S$, define $M_\theta(s):=\max_{a\in A} f_\theta(s,a)$ and $\mathcal A_\theta(s):=\operatorname*{arg\,max}_{a\in A} f_\theta(s,a)$. 
The maximum exists because $A$ is compact and $a\mapsto f_\theta(s,a)$ is continuous.
Since $a\mapsto f_\theta(s,a)-M_\theta(s)$ is a nontrivial real-analytic function on an open neighborhood of $A$, the standard zero-set theorem for real-analytic functions implies that its zero set has $d_2$-dimensional Lebesgue measure zero. 
Because $\mathcal A_\theta(s)=\{a\in A:f_\theta(s,a)-M_\theta(s)=0\}$, we obtain $\mu(\mathcal A_\theta(s))=0$ for every fixed $s\in S$ and every nonzero $\theta$.
\end{example}


\subsubsection*{Simplex features}

Assumption~\ref{asp:app:mdp:max_set_simplex} explains what we mean by simplex features and adds an additional property which allow us to get unboundedness of $\theta \mapsto \operatorname{KL}(\pi_\theta|\mu)$ in the direction orthogonal to the $\mathbf 1$-vector.
Recall that for any $v\in \mathbb R^p$ we write $v_\perp := v - p^{-1}\langle v, \mathbf 1\rangle \mathbf 1$.

\begin{restatable}{assumption}{aspSimplexMaxSet}\label{asp:app:mdp:max_set_simplex}
The feature basis $g:S\times A \to \mathbb R^p$ fall in the probability simplex $\Delta_{p-1}$ for every $s$ and $a$, i.e. $\sum_{i=1}^p g_i(s,a) = 1 $ and for all $s \in S,a \in A$, $g(s,a) \geq 0$. 
Moreover, for each fixed
$s\in S$, assume that $a\mapsto g(s,a)\in\mathbb R^p$ is continuous
and that, for every $u\in\mathbb R^p$ such that $u\neq 0$ and $u \perp \mathbf{1}$,
\[
\mu \big(\arg\max_{a\in A} u^\top g(s,a)\big)=0 .
\]
\end{restatable}
Note the assumption on the measure of the maximizer set of $u^\top g(s,a)$ is a little different between Assumption~\ref{asp:app:mdp:max_set_simplex}
and Assumption~\ref{asp:app:mdp:max_set_affine} in that we consider different vectors $u$ in each.

\begin{restatable}[Radial unboundedness of KL divergence in direction orthogonal to $\mathbf 1$]{theorem}{KLUBSIMPLEX}\label{thm:main:mdp:Radial_unboundedness_of_relative_entropy_in_subspace_simplex}
Let Assumption~\ref{asp:app:mdp:max_set_simplex} hold.
Fix $s\in S$.
Then 
\[
\lim_{\theta\in \mathbb R^p:\|\theta_\perp\|_2\to\infty}
\mathrm{KL}\bigl(\pi_\theta(\cdot\mid s)\,|\,\mu\bigr)=\infty.
\]
\end{restatable}

The proof of Theorem \ref{thm:main:mdp:Radial_unboundedness_of_relative_entropy_in_subspace_simplex} can be found in Section \ref{sec:pf_Radial_unboundedness_of_relative_entropy_in_subspace_simplex}.

Theorem \ref{thm:main:mdp:Radial_unboundedness_of_relative_entropy_in_subspace_simplex} combined with classical arguments for constructing ODE solutions with Lyapunov functions will give existence of the solution to the gradient flow~\eqref{eq:intro:continuous_GF}.

\begin{restatable}[Existence of solution to the gradient flow using simplex features]{lemma}{ExistSolSimplex}\label{lem:main:mdp:exist_sol_gf_simplex}
Let Assumption \ref{asp:app:mdp:max_set_simplex} hold.
Then there exists solution $\{\theta_t\}_{t \geq 0}$ to (\ref{eq:intro:continuous_GF}).
Moreover, $\sup_{t \geq 0} \big|\log \frac{d \pi_{\theta_t}}{d\mu}\big|_{B_b(S \times A)} < \infty$.
\end{restatable}

The proof of Lemma \ref{lem:main:mdp:exist_sol_gf_simplex} can be found in Section \ref{sec:pf_exist_sol_gf_simplex}.

\begin{restatable}{assumption}{aspInitSimplex}\label{asp:app:mdp:initial_singularity_simplex}
For all $s \in S$, $\int g(s,a)g^\top(s,a)\pi_{\theta_0} (da | s)$ is positive definite.
\end{restatable}

\begin{restatable}[Linear convergence rate using simplex features]{theorem}{LinCovSimplex}\label{thm:main:lin_cov_simplex}
Let Assumption \ref{asp:app:mdp:Q_realisability}, \ref{asp:app:mdp:max_set_simplex} and \ref{asp:app:mdp:initial_singularity_simplex} hold.
Let $\{\theta_t\}_{t \geq 0}$ be the solution to gradient flow (\ref{eq:intro:continuous_GF}).
Then there exists $C_{\theta_0} > 0$ such that 
for $C_R(\theta_t)$ from Theorem~\ref{thm:main: non-uniform_PL_ineq} we have $\sup_{t \geq 0} C_R(\theta_t) \leq C_{\theta_0}$ and
\[
0 \leq V_\tau^{\pi_{\theta_t}}(\rho) - V_\tau^{\pi^*}(\rho) \leq e^{-t C_{\theta_0}^{-1}} 
(V_\tau^{\pi_{\theta_0}}(\rho) - V_\tau^{\pi^*}(\rho)).
\]
\end{restatable}
The the exact form of $C_{\theta_0}$ can be found in Section \ref{sec:pf_lin_cov_simplex}.

To conclude discussing simplex features we provide an example that satisfies Assumption~\ref{asp:app:mdp:max_set_simplex}.

\begin{example}[Bernstein polynomial features]\label{ex:Bernstein_Poly_MDP}
Let $S\subset\mathbb R^{d_1}$ be nonempty and compact, and let $A\subset\mathbb R^{d_2}$ be compact with nonempty interior. 
Here, nonempty interior means that there exist $a_0\in A$ and $r>0$ such that the open ball $B(a_0,r)\subset A$. 
Fix $u\in\mathbb R^{d_2}\setminus\{\mathbf 0\}$ and let $q:S\to\mathbb R$ be continuous. Define $\widetilde h(s,a):=q(s)+u^\top a$.

Let $m_h:=\min_{(s,a)\in S\times A}\widetilde h(s,a)$ and $M_h:=\max_{(s,a)\in S\times A}\widetilde h(s,a)$. 
These extrema exist because $S\times A$ is compact and $\widetilde h$ is continuous. 
Moreover, $M_h>m_h$. 
Indeed, since $A$ has nonempty interior, there exist $a_0\in A$ and $r>0$ such that $B(a_0,r)\subset A$. 
Let $v:=u/\|u\|$, and define $a_+:=a_0+(r/2)v$ and $a_-:=a_0-(r/2)v$. 
Then $a_+,a_-\in A$, and for any fixed $s_0\in S$, we have $\widetilde h(s_0,a_+)-\widetilde h(s_0,a_-)=u^\top(a_+-a_-)=r\|u\|>0$. 
Hence $\widetilde h$ is not constant on $S\times A$, so its maximum and minimum are distinct. 
Define $h:S\times A\to[0,1]$ by $h(s,a):=(\widetilde h(s,a)-m_h)/(M_h-m_h)$.

Fix $k\geq 1$. 
For $\ell=0,\dots,k$, define the one-dimensional Bernstein basis function $b_{\ell,k}(w):=\binom{k}{\ell}w^\ell(1-w)^{k-\ell}$ for $w\in[0,1]$. 
Define the feature $g:S\times A\to\mathbb R^{k+1}$ by $g(s,a):=(b_{0,k}(h(s,a)),\dots,b_{k,k}(h(s,a)))^\top$. 
Then $g_\ell(s,a)\geq 0$ for every $\ell$, and $\sum_{\ell=0}^k g_\ell(s,a)=1$ for every $(s,a)\in S\times A$.

For $\theta=(\theta_0,\dots,\theta_k)\in\mathbb R^{k+1}$, define $f_\theta(s,a):=\theta^\top g(s,a)=\sum_{\ell=0}^k\theta_\ell b_{\ell,k}(h(s,a))$. 
Equivalently, $f_\theta(s,a)=P_\theta(h(s,a))$, where $P_\theta(w):=\sum_{\ell=0}^k\theta_\ell b_{\ell,k}(w)$.

Assume that $\theta\neq \mathbf 0$ and $\theta^\top\mathbf 1=\sum_{\ell=0}^k\theta_\ell=0$. 
Then $P_\theta : \mathbb [0,1] \to \mathbb R$ is non-constant. 
To see this, suppose that $P_\theta\equiv c$ on $[0,1]$. 
Since $\sum_{\ell=0}^k b_{\ell,k}(w)=1$, the constant polynomial $c$ satisfies $c=\sum_{\ell=0}^k c\,b_{\ell,k}(w)$. 
By uniqueness of expansion of the Bernstein polynomials, $\theta_\ell=c$ for every $\ell=0,\dots,k$. 
The condition $\theta^\top\mathbf 1=0$ gives $0=\sum_{\ell=0}^k\theta_\ell=(k+1)c$, so $c=0$, and hence $\theta=\mathbf0$, contradicting $\theta \neq \mathbf 0$.
Therefore $P_\theta : \mathbb [0,1] \to \mathbb R$ is non-constant.

For each fixed state $s\in S$, define $M_\theta(s):=\max_{a\in A}f_\theta(s,a)$ and $\mathcal A_\theta(s):=\operatorname*{arg\,max}_{a\in A}f_\theta(s,a)$. 
The maximum exists because $A$ is compact and $a\mapsto f_\theta(s,a)$ is continuous. 
Since $P_\theta : \mathbb [0,1] \to \mathbb R$ is a non-constant uni-variate polynomial, $P_\theta(w)-M_\theta(s)$ is not the zero polynomial. Hence the set $H_\theta(s):=\{w\in\mathbb R:P_\theta(w)=M_\theta(s)\}$ is finite.

Now, if $a\in\mathcal A_\theta(s)$, then $P_\theta(h(s,a))=M_\theta(s)$, so $h(s,a)\in H_\theta(s)$. 
Therefore $\mathcal A_\theta(s)\subseteq \bigcup_{w\in H_\theta(s)}\{a\in A:h(s,a)=w\}$. 
For fixed $s$ and $w$, the level set $\{a\in A:h(s,a)=w\}$ is contained in the affine hyperplane $\{a\in\mathbb R^{d_2}:u^\top a=m_h+(M_h-m_h)w-q(s)\}$. 
Since $u\neq  \mathbf 0$, this hyperplane has $d_2$-dimensional Lebesgue measure zero. 
Thus $\mathcal A_\theta(s)$ is contained in a finite union of measure-zero sets, and hence $\mu(\mathcal A_\theta(s))=0$.
\end{example}


\subsection*{Example of a feature basis not providing radial unboundedness} 

In bandit setting, where the MDP has a single state, we have the following example using simplex features that do not give radial unboundedness of $\operatorname{KL}$ term in the subspace orthogonal to $\mathbf 1$.
\begin{example}[Hat-functions do not give radial unboundedness in the subspace ]\label{ex:hat_func_short}
Choose the concrete grid $x_0=0,\ x_1=\frac13,\ x_2=\frac23,\ x_3=1$, and let $g_0,g_1,g_2,g_3$ be the standard one-dimensional finite-element hat functions on the grid $[0,\frac13), [\frac13,\frac23), [\frac23,1]$. Set $g(a):=(g_0(a),g_1(a),g_2(a),g_3(a))^\top$ and $\theta=(-1,1,1,-1)^\top$. Then $\theta^\top \mathbf 1=0$, and the induced function $f_\theta(a):=\theta^\top g(a)$ is
\[
f_\theta(a) = 6a-1\,\,\text{if}\, a\in\left[0,\tfrac13\right),\,\,\,
f_\theta(a) = 1\,\,\text{if}\, a\in\left[\tfrac13,\tfrac23\right),\,\,\,
f_\theta(a) = 5-6a\,\,\text{if}\, a\in\left[\tfrac23,1\right].
\]
Thus $f_\theta$ attains its maximum value $1$ on the whole interval $[\frac13,\frac23]$. Writing $\mathcal A_\theta:=\operatorname*{arg\,max}_{a\in[0,1]}f_\theta(a)$, we have $\mathcal A_\theta=[\frac13,\frac23]$ and $\mu(\mathcal A_\theta)=\frac13>0$. Hence Assumption~\ref{asp:app:mdp:max_set_simplex} fails for this choice of $\theta$.

This also shows that radial unboundedness of the entropy may fail without the zero-measure maximizer assumption. Let $\mu$ be Lebesgue measure restricted to $[0,1]$. For $\beta>0$, let $\pi_\beta$ be the probability measure whose density with respect to $\mu$ is $p_\beta(a):=\frac{d\pi_\beta}{d\mu}(a)=e^{\beta f_\theta(a)}/Z_\beta$, where $Z_\beta:=\int_0^1 e^{\beta f_\theta(a)}\,da$. 
Since $f_\theta$ is linear on the two exterior intervals and equals $1$ on $[\frac13,\frac23]$, the change of variables $y=f_\theta(a)$ gives
\[
Z_\beta
=
e^\beta
\left[
\frac13+\frac{1}{3\beta}\bigl(1-e^{-2\beta}\bigr)
\right]
=
\frac{e^\beta}{3}
\left[
1+\frac{1-e^{-2\beta}}{\beta}
\right].
\]
Moreover, we have $\mathbb E_{\pi_\beta}[f_\theta(a)]=\partial_\beta\log Z_\beta$, and therefore
\[
\operatorname{KL}(\pi_\beta\mid\mu)
=
\beta\,\mathbb E_{\pi_\beta}[f_\theta(a)]-\log Z_\beta
=
\beta\,\partial_\beta\log Z_\beta-\log Z_\beta .
\]
From the exact expression for $Z_\beta$, we have $\log Z_\beta=\beta-\log 3+\log(1+(1-e^{-2\beta})/\beta)$. Hence, as $\beta\to\infty$,
\[
\log Z_\beta
=
\beta-\log 3+\frac1\beta+\mathcal O(\beta^{-2}),
\qquad
\beta\,\partial_\beta\log Z_\beta
=
\beta-\frac1\beta+\mathcal O(\beta^{-2}).
\]
Consequently, $\operatorname{KL}(\pi_\beta\mid\mu) = \log 3-\frac{2}{\beta}+\mathcal O(\beta^{-2})$,
and in particular $\lim_{\beta\to\infty}\operatorname{KL}(\pi_\beta\mid\mu)=\log 3$, which is finite.
\end{example}
The detailed calculation can be found in Section \ref{sec:ex}.








\allowdisplaybreaks
\section{Proofs}
\subsection{Basic notations and definitions}
\label{sec:nota and def}

 
For matrix $G \in \mathbb{R}^{p \times p}$, denote $\lambda_{\min} (G)$ the smallest eigenvalue value of $G$. 
Let $\mathbb S^{p-1}
:=
\left\{u\in\mathbb R^p:\|u\|_2=1\right\}.$ 
Let $ \mathbf{1} = (1,1, \dots,1) \in \mathbb{R}^p$. For vectors $\mathbf{a}, \mathbf{b} \in \mathbb{R}^p$, the inner product $\langle \mathbf{a}, \mathbf{b}\rangle$ of $\mathbf{a}\ \text{and} \ \mathbf{b}$ is $\mathbf{a}^\top\mathbf{b}$.
And we use both notations throughout the paper.

Let $(E,d)$ denote a complete separable metric space (i.e. a Polish space). 
For a given measure $\rho$ in $E$, denote by $L^p(E,\rho)$, $p\in [1,\infty]$, for Lebesgue spaces of integrable functions. 
We always equip a Polish space with its Borel sigma-field $\mathcal{B}(E).$ 
Denote by $B_b(E)$ the space of bounded strongly measurable functions $f:E\rightarrow \mathbb{R}$ endowed with the supremum norm $|f|_{B_b(E)}=\sup_{x\in E}|f(x)|$. 
Denote by $\mathcal{M}(E)$ the Banach space of signed measures (finite) $\mu$ on $E$ endowed with the total variation norm $|\mu|_{\mathcal{M}(A)}=|\mu|(E)$, where $|\mu|$ is the total-variation measure. 
We note that if $\mu=f d\rho$, where $\rho \in \mathcal{M}_+(E)$ is a non-negative measure and $f\in L^1(E,\rho)$, then  $|\mu|_{\mathcal{M}(E)}=|f|_{L^1(E,\rho)}$. 
We denote by $\mathcal{P}(E)\subset\mathcal{M}(E)$ the convex subset of probability measures on $E$. 
For $\mu,\mu'\in \mathcal{P}(E)$ such that $\mu$ is absolutely continuous with respect to $\mu'$, the relative entropy of $\mu$ with respect to $\mu'$ (or KL divergence of $\mu$ relative to $\mu'$) is defined by 
\[
\textnormal{KL}(\mu|\mu')=\int_{E}\log\frac{d\mu}{d\mu'}(x)\mu(dx)\,.
\]

It is convenient to have notation for measurable functions $k:E_1\rightarrow \mathcal{M}(E_2)$ for given Polish spaces $(E_1,d_1)$ and $(E_2,d_2)$. 
For example, $P:S\rightarrow \mathcal{P}(S\times A)$ will denote a controlled transition probability and $\pi: S\rightarrow \mathcal{P}(A)$ a stochastic policy. 
Denote by $b\mathcal{K}(E_1|E_2)$ the Banach space of bounded signed kernels $k: E_2 \rightarrow \mathcal{M}(E_1)$ endowed with the norm $|k|_{b\mathcal{K}(E_1|E_2)}=\sup_{x\in E_2}|k(x)|_{\mathcal{M}(E_1)}$; that is, $k(U|\cdot): E_2\rightarrow \mathbb{R}$ is measurable for all $U\in \mathcal{M}(E_1)$ and $k(\cdot|x)\in \mathcal{M}(E_1)$ for all $x\in E_2$. 
For a fixed positive reference measure $\mu \in \mathcal{M}(E_1)$, we denote by $b\mathcal{K}_{\mu}(E_1|E_2)$ the space of bounded kernels that are absolutely continuous with respect to $\mu$. 

Every kernel $k\in b\mathcal{K}(E_1|E_2)$ induces bounded linear operators $T_k \in \mathcal{L}(\mathcal{M}(E_2),\mathcal{M}(E_1))$ and $S_k \in \mathcal{L}(B_b(E_1),B_b(E_2))$ defined by
\[
T_k\mu(dy)=\mu k(dy)=\int_{E_2}\mu(dx)k(dy|x)
\]
and 
\[
S_kf(x)=\int_{E_1}k(dy|x)f(y)\,,
\]
respectively.
Moreover, by Exercise 2.3 and Proposition 3.1 in~\cite{kunze2011pettis}, we have
\begin{align*}
|k|_{b\mathcal{K}(E_1|E_2)}&=\sup_{x\in E_2}\underset{|h|_{B_b(E_1)}\le 1}{\sup_{h\in B_b(E_1)}}\int_{E_1}h(y)k(dy|x)\\
&=|S_k|_{\mathcal{L}(B_b(E_1),B_b(E_2))}\\
&=|T_k|_{\mathcal{L}(\mathcal{M}(E_2),\mathcal{M}(E_1))}\,,
\end{align*}
where the latter are operator norms. 
Thus, $b\mathcal{K}(E|E)$ is a Banach algebra with the product defined via composition of the corresponding linear operators; in particular, for a given $k\in b\mathcal{K}(E|E)$,
\[
\begin{split}
& T_k^n\mu(dy)=\mu k^n(dy)\\
& = \int_{E^{n}}\mu(dx_0)k(dx_1|x_0)\cdots k(dx_{n-1}|x_{n-2}) k(dy|x_{n-1})\,.
\end{split}
\]
Notice that if $f\in L^\infty(E_1,\mu)$ and $k\in b\mathcal{K}_{\mu}(E_1|E_2)$, then for all $x\in E_2$,
\begin{align}
S_kf(x)&=\int_{E_1}\mu(dy)\frac{dk}{d\mu}(y|x)f(y)\notag\\
&\le |f|_{L^\infty(E_1,\mu)} \left|\frac{dk}{d\mu}(\cdot|x)\right|_{L^1(E_1,\mu)}\notag \\
&\le |f|_{L^\infty(E_1,\mu)} |k|_{b\mathcal{K}(E_1|E_2)}\,. \label{ineq:app:essup_kernel}
\end{align}

We denote by $\mathcal{P}(E_1|E_2)$ the convex subspace of $P\in b\mathcal{K}(E_1|E_2)$ such that $P(\cdot|x)\in \mathcal{P}(E_1)$ for all $x\in E_2$; such kernels are referred to as stochastic kernels. 
A stochastic kernel $P\in \mathcal{P}(E_1|E_2)$ is said to be strongly Feller if $\int_{E_1}P(dy|x)f(y)$ is continuous in $x\in E_2$ for all $f\in B_b(E_1)$. For a fixed positive reference measure $\mu \in \mathcal{M}(E_1)$, we denote by $\mathcal{P}_{\mu}(E_1|E_2)$ the space of kernels that are absolutely continuous with respect to $\mu$. 
A bounded kernel $k\in b\mathcal{K}(E_1|E_2)$ is thus strongly Feller if the range of $S_k$ lies in the space of continuous functions on $E_2$. 

\subsection{Basic results on entropy regularized MDPs}

The following lemma (see e.g., Lemma 2.3 \cite{kerimkulov2025fisher}) is crucial for addressing this non-convexity issue. 
\begin{lemma}[Performance difference]
\label{lem:app:mdp:performance_diff}
For all $\rho \in \mathcal{P}(S)$ and $\pi,\pi'\in \Pi_{\mu}$, 
\begin{align*}
&V^{\pi}_\tau(\rho)-V^{\pi'}_\tau(\rho) \\
& = \frac{1}{1-\gamma}\int_S \bigg[\int_A\left(Q^{\pi'}_{\tau}(s,a)+\tau \log \frac{d \pi'}{d\mu}(a|s)\right)(\pi-\pi')(da|s) + \tau    \operatorname{KL}(\pi(\cdot | s)|\pi'(\cdot | s)) \bigg]d^{\pi}_\rho(ds)\,.
\end{align*}
\end{lemma}

Define the proximal policy
\begin{equation}\label{eq:app:mdp:proximal_policy}
\pi_{\pi^{\prime}}(d a | s)=\frac{1}{Z_{\pi^{\prime}}(s)} \exp \left(-\frac{1}{\tau}\left(Q^{\pi^{\prime}}_\tau(s, a)-V^{\pi^{\prime}}_
\tau(s)\right)\right) \mu(d a),
\end{equation}
where $Z_{\pi^{\prime}}(s):=\int_A \exp \left(-\frac{1}{\tau}\left(Q^{\pi^{\prime}}_\tau\left(s, a^{\prime}\right)-V^{\pi^{\prime}}_\tau(s)\right)\right) \mu\left(d a^{\prime}\right)$.
Then we have the following sandwich inequality of sub-optimal gap in terms of KL divergence:
\begin{lemma}[Sandwich Inequality]\label{lem:sandwich-inequality}
Let $\pi_{\bar\pi}$ denote the proximal policy step associated with a policy
$\bar\pi$. 
Then, for any policies $\pi,\pi'$ and any initial distribution
$\rho$,
\begin{equation}
V_\tau^{\pi'}(\rho)-V_\tau^\pi(\rho)
\le
\frac{\tau}{1-\gamma}
\int_S
\mathrm{KL}(\pi'(\cdot | s)| \pi_{\pi'}(\cdot |s))
\,d_\rho^\pi(ds).
\label{eq:sandwich-general-bound}
\end{equation}
Moreover, if $\pi^*$ is an optimal policy, then for any policy $\pi'$ and any
initial distribution $\rho$,
\begin{equation}
\frac{\tau}{1-\gamma}
\int_S
\mathrm{KL}(\pi'(\cdot | s)| \pi^*(\cdot | s))
\,d_\rho^{\pi'}(ds)
=
V_\tau^{\pi'}(\rho)-V_\tau^{\pi^*}(\rho)
\le
\frac{\tau}{1-\gamma}
\int_S
\mathrm{KL}(\pi'(\cdot | s)| \pi_{\pi'}(\cdot |s))
\,d_\rho^{\pi^*}(ds).
\label{eq:sandwich-inequality}
\end{equation}
\end{lemma}

\begin{proof}
The flat derivative of the objective can be written in terms of the proximal
policy (\ref{eq:app:mdp:proximal_policy}) as
\begin{align}
\frac{\delta V_\tau^{\pi'}}{\delta \pi}
&=
Q_\tau^{\pi'}-V_\tau^{\pi'}
+
\tau \log \frac{d\pi'}{d\mu}
\nonumber \\
&=
\tau \log \frac{d\pi'}{d\mu}
-
\tau
\log
\exp\left(
-\frac{1}{\tau}
\left(Q_\tau^{\pi'}-V_\tau^{\pi'}\right)
\right)
+
\tau \log Z_{\pi'}
-
\tau \log Z_{\pi'}
\nonumber \\
&=
\tau \log \frac{d\pi'}{d\mu}
-
\tau \log \frac{d\pi_{\pi'}}{d\mu}
-
\tau \log Z_{\pi'}
=
\tau \log \frac{d\pi'}{d\pi_{\pi'}}
-
\tau \log Z_{\pi'} .
\label{eq:sandwich-flat-derivative}
\end{align}

By the regularized performance-difference Lemma \ref{lem:app:mdp:performance_diff}, for any $\pi,\pi' \in \Pi_\mu$ and
any $\rho$,
\begin{align}
V_\tau^\pi(\rho)
&=
V_\tau^{\pi'}(\rho)
+
\frac{1}{1-\gamma}
\int_S
\left[
\int_A
\frac{\delta V_\tau^{\pi'}}{\delta \pi}(s,a)
(\pi-\pi')(da| s)
+
\tau \mathrm{KL}(\pi(\cdot | s)| \pi'(\cdot | s))
\right]
d_\rho^\pi(ds).
\label{eq:sandwich-performance-difference}
\end{align}
Using the expression of the flat derivative, we obtain
\begin{align}
V_\tau^\pi(\rho)
&=
V_\tau^{\pi'}(\rho)
+
\frac{1}{1-\gamma}
\int_S
\bigg[
\int_A
\left(
Q_\tau^{\pi'}(s,a)
+
\tau \log \frac{d\pi'}{d\mu}(a|s)
-
V_\tau^{\pi'}(s)
\right)
(\pi-\pi')(da| s)
\nonumber \\
&\hspace{4cm}
+
\tau \mathrm{KL}(\pi(\cdot | s)| \pi'(\cdot |s))
\bigg]
d_\rho^\pi(ds)
\nonumber \\
&=
V_\tau^{\pi'}(\rho)
+
\frac{1}{1-\gamma}
\int_S
\bigg[
\int_A
\left(
Q_\tau^{\pi'}(s,a)
+
\tau \log \frac{d\pi'}{d\mu}(a|s)
\right)
(\pi-\pi')(da| s)
\nonumber \\
&\hspace{4cm}
+
\tau \mathrm{KL}(\pi(\cdot | s)| \pi'(\cdot |s))
\bigg]
d_\rho^\pi(ds)
\nonumber \\
&=
V_\tau^{\pi'}(\rho)
+
\frac{1}{1-\gamma}
\int_S
\bigg[
\int_A
\left(
Q_\tau^{\pi'}(s,a)
+
\tau \log \frac{d\pi'}{d\mu}(a | s)
\right)
(\pi-\pi')(da| s)
\nonumber \\
&\hspace{4cm}
+
\int_A
\tau \log \frac{d\pi}{d\pi'}(a | s)
\,\pi(da| s)
\bigg]
d_\rho^\pi(ds)
\nonumber \\
&=
V_\tau^{\pi'}(\rho)
+
\frac{1}{1-\gamma}
\int_S
\bigg[
\int_A
Q_\tau^{\pi'}(s,a)
(\pi-\pi')(da| s)
+
\tau \mathrm{KL}(\pi(\cdot |s)| \mu)
\nonumber \\
&\hspace{4cm}
-
\tau
\int_A
\log \frac{d\pi'}{d\mu}(a|s)
\,\pi'(da| s)
\bigg]
d_\rho^\pi(ds).
\label{eq:sandwich-expanded-difference}
\end{align}

Minimizing the term inside the integral over $\pi$ for every $s\in S$ and
recalling the proximal policy step gives
\begin{align}
V_\tau^\pi(\rho)
&\ge
V_\tau^{\pi'}(\rho)
+
\frac{1}{1-\gamma}
\int_S
\left[
\int_A
\frac{\delta V_\tau^{\pi'}}{\delta \pi}(s,a)
(\pi_{\pi'}-\pi')(da| s)
+
\tau \mathrm{KL}(\pi_{\pi'}(\cdot | s)| \pi'(\cdot | s))
\right]
d_\rho^\pi(ds).
\label{eq:sandwich-proximal-lower-bound}
\end{align}
From \eqref{eq:sandwich-flat-derivative}, for any $\pi,\pi'$,
\begin{align}
&\int_A
\frac{\delta V_\tau^{\pi'}}{\delta \pi}(s,a)
(\pi_{\pi'}-\pi')(da| s)
+
\tau \mathrm{KL}(\pi_{\pi'}(\cdot | s)| \pi'(\cdot | s))
\nonumber \\
&=
\tau
\int_A
\log \frac{d\pi'}{d\pi_{\pi'}}(a|s)
(\pi_{\pi'}-\pi')(da| s)
+
\tau
\int_A
\log \frac{d\pi_{\pi'}}{d\pi'}(a|s)
\,\pi_{\pi'}(da| s)
\nonumber \\
&=
-\tau
\int_A
\log \frac{d\pi_{\pi'}}{d\pi'}(a|s)
(\pi_{\pi'}-\pi')(da| s)
+
\tau
\int_A
\log \frac{d\pi_{\pi'}}{d\pi'}(a|s)
\,\pi_{\pi'}(da| s)
\nonumber \\
&=
\tau
\int_A
\log \frac{d\pi_{\pi'}}{d\pi'}(a|s)
\,\pi'(da| s)
=
-\tau \mathrm{KL}(\pi'(\cdot | s)| \pi_{\pi'}(\cdot |s)).
\label{eq:sandwich-kl-cancellation}
\end{align}
Combining \eqref{eq:sandwich-proximal-lower-bound} and
\eqref{eq:sandwich-kl-cancellation}, we get
\begin{equation}
V_\tau^\pi(\rho)
\ge
V_\tau^{\pi'}(\rho)
-
\frac{\tau}{1-\gamma}
\int_S
\mathrm{KL}(\pi'(\cdot | s)| \pi_{\pi'}(\cdot |s))
\,d_\rho^\pi(ds).
\label{eq:sandwich-lower-bound}
\end{equation}
Equivalently,
\begin{equation}
V_\tau^{\pi'}(\rho)-V_\tau^\pi(\rho)
\le
\frac{\tau}{1-\gamma}
\int_S
\mathrm{KL}(\pi'(\cdot | s)| \pi_{\pi'}(\cdot |s))
\,d_\rho^\pi(ds),
\label{eq:sandwich-value-gap-upper}
\end{equation}
which proves \eqref{eq:sandwich-general-bound}.

Finally, the optimality condition for $\pi^*$ gives, for all $s \in S,a \in A$,
\begin{equation}
\frac{\delta V_\tau^{\pi^*}}{\delta \pi}
=
Q_\tau^{\pi^*}
-
V_\tau^{\pi^*}
+
\tau \log \frac{d\pi^*}{d\mu}
=
0.
\label{eq:sandwich-optimality-condition}
\end{equation}
Using this in the regularized performance-difference identity yields
\begin{equation}
V_\tau^{\pi'}(\rho)-V_\tau^{\pi^*}(\rho)
=
\frac{\tau}{1-\gamma}
\int_S
\mathrm{KL}(\pi'(\cdot |s)| \pi^*(\cdot | s))
\,d_\rho^{\pi'}(ds).
\label{eq:sandwich-optimal-gap}
\end{equation}
Taking $\pi=\pi^*$ in \eqref{eq:sandwich-value-gap-upper} and combining with
\eqref{eq:sandwich-optimal-gap} gives
\eqref{eq:sandwich-inequality}.
\end{proof}



\begin{lemma}[Gradient of log policy]\label{lem:app:mdp:grad_of_log_policy}
For any $\theta \in \mathbb{R}^p$, $s \in S$, $a \in A$
\begin{equation}\label{eq:app:mdp:grad_of_log_policy}
\nabla \log\frac{d \pi_{\theta}}{d \mu}(a|s) = g(s,a) - \int_A g(s, a^\prime) \pi_{\theta}( da^\prime | s). 
\end{equation}
\end{lemma}
\begin{proof}
For any  $s \in S$, $a \in A$,
\[
\log\frac{d \pi_{\theta}}{d \mu}(a|s) = \langle\theta,g(s,a)\rangle - \log \int_A \exp\left(\langle\theta,g(s,a)\rangle\right)\mu(da),
\]
then taking gradient concludes the proof.
\end{proof}
In the remainder of the paper, $\nabla$ means $\nabla_\theta$ unless stated otherwise.
\begin{definition}\label{def:app:mdp:FIM}
For fixed $s$ and for $\theta \in \mathbb{R}^p$, the Fisher information matrix (FIM) $G^{\pi_\theta}(s)$ is defined as
\begin{align*}
G^{\pi_\theta}(s) & = \int_A \left(\nabla \log\frac{d \pi_{\theta}}{d \mu}(a|s) \right)\left(\nabla \log\frac{d \pi_{\theta}}{d \mu}(a|s)\right)^{\top}\pi_\theta(da|s)\\
& = \int_A \left(g(s,a) - \int_A g(s, a^\prime) \pi_\theta ( da^\prime | s) \right)\left(g(s,a) - \int_A g(s, a^\prime) \pi_\theta (da^\prime | s) \right)^{\top}\pi_\theta(da|s),
\end{align*}
and for $\rho \in \mathcal{P}(S)$,$G^{\pi_\theta}(\rho) = \int G^{\pi_\theta}(s) d_\rho^{\pi_\theta}(ds)$. 
\end{definition}

\begin{lemma}\label{lem:app:mdp:bounded_grad_of_log_policy}
For any $\theta \in \mathbb{R}^p$ and any $s \in S$ and $a \in A$,
\begin{equation}
\left\|\nabla \log\frac{d \pi_{\theta}}{d \mu} (a|s)\right\|_2 \leq 2.
\end{equation}
\end{lemma}

\begin{lemma}[Smoothness of $\theta \mapsto \nabla \log\frac{d \pi_{\theta}}{d \mu} (a|s)$]\label{lem: smooth of log policy}
For any $\theta, \theta^\prime \in \mathbb{R}^p$, then for any $s \in S$ and $a \in A$,        \begin{equation}\label{eq: smooth of log policy}
\left\|\nabla \log\frac{d \pi_{\theta^\prime}}{d \mu} (a|s) - \nabla \log\frac{d \pi_{\theta}}{d \mu} (a|s)\right\|_2 \leq 2\|\theta^\prime - \theta\|_2. 
\end{equation}
\end{lemma}
\begin{proof}
\begin{align*}
& \nabla \log\frac{d \pi_{\theta^\prime}}{d \mu} (a|s) - \nabla \log\frac{d \pi_{\theta}}{d \mu} (a|s) \\
& = \int_A g(s,a) \left(\pi_{\theta^\prime} - \pi_{\theta}\right)(da|s)\\
& =  \int_0^1 \int_A g(s,a) \bigg\langle\nabla \log\frac{d \pi_{\theta^\epsilon}}{d \mu} (a|s),\theta^\prime - \theta \bigg\rangle \pi_{\theta^\epsilon} (da|s)d \epsilon  \\
& = \int_0^1 \int_A g(s,a) \left\langle g(s, a) - \int_A g(s,a^\prime)\pi_{\theta^\epsilon} (da|s),\theta^\prime - \theta \right\rangle \pi_{\theta^\epsilon} (da|s)d \epsilon \\
& = \int_0^1 \int_A \langle g(s,a) , g(s,a) - \int_A g(s,a^\prime)\pi_{\theta^\epsilon} (da|s)\rangle \pi_{\theta^\epsilon} (da|s)d \epsilon \cdot (\theta^\prime - \theta),
\end{align*}
where we use Lemma \ref{lem:app:mdp:grad_of_log_policy} in the third equality.
Hence by the bound of $g$,
\begin{align*}
&\left\|\nabla \log\frac{d \pi_{\theta^\prime}}{d \mu} (s,a) - \nabla \log\frac{d \pi_{\theta}}{d \mu} (a|s)\right\|_2\\
&\leq \left|\sup_a \|g(s,a)\|^2_2 - \left\|\int_0^1 \int_A g(s,a)\pi_{\theta^\epsilon} (da|s)d \epsilon\right\|^2_2\right| \cdot \|\theta^\prime - \theta\|_2\\
& \leq 2\|\theta^\prime - \theta\|_2,
\end{align*}
which concludes the proof.
\end{proof}

\begin{lemma}\label{lem:bounds}
For any $\theta \in \mathbb{R}^p$, we have
\begin{gather*}
\left|V^{\pi_\theta}_{\tau}(s)\right|\le \frac{1}{1-\gamma} \left(1+\tau \bigg|\log \frac{d\pi_{\theta}}{d \mu}\bigg|_{B_b(S \times A)} \right),\\
|Q^{\pi_{\theta}}_{\tau}(s,a)| \le \frac{1}{1-\gamma} \left(1+\gamma\tau \bigg|\log \frac{d\pi_{\theta}}{d \mu}\bigg|_{B_b(S \times A)} \right).
\end{gather*}
\end{lemma}
\begin{proof}

By the bound on cost function and Eq.~\eqref{eq:intro:value_function_occ}, we obtain
\begin{equation}\label{eq:bound_on_V_pi}
\begin{aligned}
|V^{\pi_\theta}_{\tau}(s)|&=\frac{1}{1-\gamma}\left|\int_{S}\int_A \left(c(s',a)+\tau \log \frac{d\pi_{\theta}}{d\mu}(a|s)\right)\pi_{\theta}(da|s')d^{\pi_\theta}(ds'|s)\right|\\
&\le \frac{1}{1-\gamma} \left(1+\tau \bigg|\log \frac{d\pi_{\theta}}{d \mu}\bigg|_{B_b(S \times A)} \right)\,.
\end{aligned}
\end{equation}
To estimate the state-action value function, by Eq.~\eqref{def:intro:Qpi}, we have
\begin{align*}
|Q^{\pi_{\theta}}_{\tau}(s,a)| &\le \frac{1}{1-\gamma} \left(1+\gamma\tau \bigg|\log \frac{d\pi_{\theta}}{d \mu}\bigg|_{B_b(S \times A)} \right)\,.
\end{align*}
\end{proof}

Next we introduce the following Lipschitz continuity of the occupancy kernel (See Lemma A.4 in~\cite{leahy2022convergence}). 
\begin{lemma}\label{lem:app:mdp:lip_occ_kernel}
For given $\pi,\pi' \in \mathcal{P}(A|S)$, we have
\[
|d^{\pi'}-d^{\pi}|_{b\mathcal{K}(S|S)}\le \frac{\gamma}{1-\gamma}|\pi'-\pi|_{b\mathcal{K}(A|S)}\,.
\]
\end{lemma}

\begin{corollary}[Lipschitz continuity of the occupancy measure in the parameter]\label{cor:app:mdp:Lip_occ}
For given $\theta,\theta' \in \mathbb{R}^p$, we have
\[
|d^{\pi_{\theta'}}-d^{\pi_{\theta}}|_{b\mathcal{K}(S|S)}\le \frac{2\gamma}{1-\gamma}\|\theta' - \theta\|_2\,.
\]
\end{corollary}
\begin{proof}
By Lemma \ref{lem:app:mdp:lip_occ_kernel}, it is enough to show that
\begin{equation}\label{ineq:Lip_pi}
|\pi_{\theta'}-\pi_{\theta}|_{b\mathcal{K}(A|S)} \le 2\|\theta' - \theta\|_2\,.
\end{equation}
We have 
\begin{align*}
|\pi_{\theta'}-\pi_{\theta}|_{b\mathcal{K}(A|S)}&=\sup_{s\in S}|\pi_{\theta'}(\cdot|s)-\pi_{\theta}(\cdot|s)|_{\mathcal{M}(A)}\\
&=\sup_{s\in S}\underset{|h|_{B_b(A)}\le 1}{\sup_{h\in B_b(A)}}\int_{A}h(a)\left(\pi_{\theta'}-\pi_{\theta}\right)(da|s)\,.
\end{align*}
Let $\theta^{\varepsilon}=\varepsilon \theta'+(1-\varepsilon)\theta$, $\varepsilon\in [0,1]$. Let $s\in S$ and $h\in B_b(A)$ be arbitrarily given. 
Using Lemma \ref{lem:app:mdp:grad_of_log_policy}, we find that 
\[
\int_{A}h(a)\left(\pi_{\theta'}-\pi_{\theta}\right)(da|s)= \left\langle\int_{A} h(a)\int_0^1 \nabla \log \frac{d \pi_{\theta^\varepsilon}}{d \mu} (a|s) \pi_{\theta^\varepsilon}(da|s)d\varepsilon, \theta^\prime - \theta\right\rangle,
\]
which results into (\ref{ineq:Lip_pi}) by Lemma \ref{lem:app:mdp:bounded_grad_of_log_policy} and concludes the proof.
\end{proof}

The proof of Proposition~\ref{prp:app:mdp:Gradient_of_the_Objective} follows Lemma 2.3 in~\cite{leahy2022convergence} and can be found in Section \ref{sec:add_proof}.
It verifies that chain rule holds for $\pi \mapsto V^\pi_\tau$ and $\theta \mapsto \pi_\theta$.

\begin{proposition}[Gradient of the Objective]\label{prp:app:mdp:Gradient_of_the_Objective} 
For any $\theta \in \mathbb{R}^p$, 
\begin{equation}\label{eq:wts_func_deriv_V}
\nabla V_\tau^{\pi_\theta}(\rho) = \frac{1}{1-\gamma}\int_S\int_A \left(Q^{\pi_\theta}_\tau(s,a) + \tau \log \frac{d \pi_{\theta}}{d \mu}(a|s)\right)\nabla \log \frac{d\pi_\theta}{d \mu}(a| s)\pi_{\theta}(da|s)d_\rho^{\pi_{\theta}}(ds).
\end{equation}
\end{proposition}

\begin{proposition}[Local Lipschitz Continuity of  $\theta \mapsto \nabla V_\tau^{\pi_\theta}(\rho)$]\label{prp:app:mdp:local_lipchitiz_conitinuity_mdp}
Let $R >0$.
For any $\theta \in \mathbb R^p$ such that $\left|\log \frac{d \pi_\theta}{d\mu}\right|_{B_b(S \times A)} \leq R$, $\theta \mapsto \nabla V_\tau^{\pi_\theta}(\rho)$ is Lipschitz continuous:
\begin{align*}
\left\|\nabla V_\tau^{\pi_\theta}(\rho) - \nabla V_\tau^{\pi_{\theta^\prime}}(\rho)\right\|_2 
\leq  C_{\gamma,\tau,R} \left\|\theta - \theta^\prime\right\|_2,
\end{align*}
where $C_{\gamma,\tau,R} = \left( \frac{1}{1-\gamma}\left(\frac{ \gamma (5 + \tau R)  }{1-\gamma} + 6\right)\left(\frac{1+\gamma\tau R}{1-\gamma} + \tau R\right) + \frac{2\tau }{1-\gamma}\right)$.
\end{proposition}
The proof of Proposition~\ref{prp:app:mdp:local_lipchitiz_conitinuity_mdp} can be found in Section~\ref{sec:proof_local_lip}

\subsection{Proof of the main results}
\label{sec:proof}

\subsubsection{Proofs of Theorem \ref{thm:main: non-uniform_PL_ineq}}
\label{sec:pf_non_uniform_PL}
Let us restate Theorem \ref{thm:main: non-uniform_PL_ineq} in its full version.


\textbf{Theorem \ref{thm:main: non-uniform_PL_ineq}.}
Let Assumption~\ref{asp:app:mdp:Q_realisability} hold.
Let $R \geq |\log \frac{d \pi_{\theta^\ast}}{d\mu}|_{B_b(S\times A)}$.
Let
\[
\Theta_R = \{\theta: \pi_\theta \in \Pi_\mu \ \text{and}\  \vert \log \frac{d \pi_\theta}{d\mu}\vert_{B_b(S\times A)} \leq R \}\,.
\]
Then for any $\theta \in \Theta_R$ there exists $C_R(\theta)>0$ such that
\begin{equation}
\label{eq:pl_theorem_appendix}
0 \leq V_\tau^{\pi_\theta}(\rho) - V_\tau^{\pi_{\theta^*}}(\rho) \leq C_R(\theta)\left\|\nabla V_\tau^{\pi_\theta}(\rho)\right\|^2_2,
\end{equation}
where 
\[
C_R(\theta) = \frac{1}{2(1-\gamma)\tau \exp\left({-\max\left(\frac{2}{(1-\gamma)\tau} \left(1+\gamma\tau R\right), R\right)}\right)}\left|\frac{d d_\rho^{\pi_{\theta^*}}}{ d \rho}\right|^2_{B_b(S)} \lambda^{\theta}\, , \ \label{eq:non_uni_PL_cons_affine}
\]
and
$
\lambda^{\theta} = \bigg[\int_S \lambda_{\min}(G^{\pi_\theta}(s))d_\rho^{\pi_{\theta^*}}(ds)\bigg]^{-2}.
$
Moreover, if the feature basis $g:S\times A \to \mathbb R^p$ fall in the probability simplex $\Delta_{p-1}$ for every $s$ and $a$, i.e. $\sum_{i=1}^p g_i(s,a) = 1 $ and for all $s \in S,a \in A$, $g(s,a) \geq 0$, then we may replace $\lambda^\theta$ in~\eqref{eq:pl_theorem_appendix} with
$\lambda^{\theta} = \left[\int_{S}\lambda_{\min}\left(\int_{A} g(s,a)g^\top(s,a)\pi_\theta(da|s)\right)d^{\pi_{\theta}}_{\rho}(ds)\right]^{-2}$.



\begin{proof}
Define proximal policy $\bar{\pi}_\theta(da|s) \propto \exp\left(-\frac{1}{\tau}Q^{\pi_\theta}_\tau(s,a) \right)\mu(da)$ as~\eqref{eq:app:mdp:proximal_policy} where by Lemma \ref{lem:bounds}, we have
\[
|Q^{\pi_{\theta}}_{\tau}(s,a)| \le \tfrac{1}{1-\gamma} \left(1+\gamma\tau R\right),
\]
then we have
\[
\big|\log \tfrac{d\bar{\pi}_{{\theta}}}{d\mu}(a|s)\big| \leq \tfrac{2}{(1-\gamma)\tau} \left(1+\gamma\tau R\right).
\]
Also, from Assumption \ref{asp:app:mdp:Q_realisability}, we know there exists $\boldsymbol{\theta}(\pi_\theta)$ such that $\langle \boldsymbol{\theta}(\pi_\theta), g(s,a)\rangle = -\frac{1}{\tau} Q^{\pi_\theta}_\tau(s,a)$. 

Hence by Lemma \ref{lem:sandwich-inequality} and Lemma \ref{lem:app:kl-logit_ineq_affine}, we have
\begin{align}
\label{eq:sandwich_ineq}
V_\tau^{\pi_\theta}(\rho) - V_\tau^{\pi_{\theta^*}}(\rho) 
\leq \frac{\tau}{1-\gamma}\int_S \mathrm{KL}(\pi_\theta (\cdot |s) | \bar{\pi}_{\theta}(\cdot |s))d_\rho^{\pi_{\theta^*}}(ds) 
& \leq \frac{\tau e^{\hat{R}} }{2(1-\gamma) }\left\|\theta - \boldsymbol{\theta}(\pi_\theta) \right\|_2^2\, , 
\end{align}
where $\hat{R} = \max\left(\frac{2}{(1-\gamma)\tau} \left(1+\gamma\tau R\right), R\right)$.
On the other hand, by Proposition \ref{prp:app:mdp:Gradient_of_the_Objective},
\begin{align}
& \nabla V_\tau^{\pi_\theta}(\rho) \nonumber \\
& = \frac{1}{1-\gamma}\int_S\int_A \left(Q^{\pi_\theta}_\tau(s,a) + \tau \log \frac{d \pi_{\theta}}{d \mu}(a|s)\right)\nabla \log \frac{d\pi_\theta}{d \mu}(a| s)\pi_{\theta}(da|s)d_{\rho}^{\pi_{\theta}}(s) \nonumber \\
& = \frac{\tau}{1-\gamma}\int_S\int_A \left(\left \langle \theta - \boldsymbol{\theta}(\pi_\theta),g(s,a)\right  \rangle - Z_{\pi_\theta}\right)\nabla \log \frac{d\pi_\theta}{d \mu}(a| s)\pi_{\theta}(da|s)d_{\rho}^{\pi_{\theta}}(s) \nonumber \\
& = \frac{\tau}{1-\gamma}\int_S\int_A \left \langle \theta - \boldsymbol{\theta}(\pi_\theta),g(s,a)- \int g(s,a^\prime) \pi_\theta(da^\prime |s)\right  \rangle  \nabla \log \frac{d\pi_\theta}{d \mu}(a| s)\pi_{\theta}(da|s)d_{\rho}^{\pi_{\theta}}(s) \nonumber \\
& = \frac{\tau}{1-\gamma}\int_S\int_A \left \langle \theta - \boldsymbol{\theta}(\pi_\theta),\nabla \log \frac{d\pi_\theta}{d \mu}(a| s)\right  \rangle \nabla \log \frac{d\pi_\theta}{d \mu}(a| s)\pi_{\theta}(da|s)d_{\rho}^{\pi_{\theta}}(s) \nonumber \\
& = \frac{\tau}{1-\gamma} \int_{S} G^{\pi_\theta}(s)d_{\rho}^{\pi_{\theta}}(ds)\cdot(\theta -\boldsymbol{\theta}(\pi_\theta))\, , \ \label{eq:grad_for_PL}
\end{align}
where the second and third equality above is due to Lemma \ref{lem:app:mdp:grad_of_log_policy} and the fact that $\int_A (g(s,a) - \int_A g(s,a^\prime)\pi_\theta(da^\prime))\pi_\theta(da) = \mathbf 0$. 
Left multiplying the above identity by $(\theta-\boldsymbol{\theta}(\pi_\theta))^\top$, using that the FIM is positive semi-definite, using the Cauchy--Schwartz inequality and finally dividing by $\|\theta-\boldsymbol{\theta}(\pi_\theta)\|_2$ we have
\begin{align*}
\left\|\nabla V_\tau^{\pi_\theta}(\rho)\right\|_2 & \geq \frac{\tau}{1-\gamma} \left|\frac{d d_\rho^{\pi_{\theta^*}}}{d d_{\rho}^{\pi_{\theta}}}\right|_{B_b(S)}^{-1} \int_{S} \lambda_{\min}\left(G^{\pi_\theta}(s)\right)d_{\rho}^{\pi_{\theta^*}}(ds)\left\|\theta -\boldsymbol{\theta}(\pi_\theta)\right\|_{2} .
\end{align*}
Due to~\eqref{eq:sandwich_ineq} and this we get
\begin{align*}
& V_\tau^{\pi_\theta}(\rho) - V_\tau^{\pi_{\theta^*}}(\rho)\\
& \leq \frac{\tau }{2(1-\gamma)e^{-\hat{R}}}\left\|\theta - \boldsymbol{\theta}(\pi_\theta)\right\|_{2}^2\\
& \leq \frac{1-\gamma}{2\tau e^{-\hat{R}}}\left|\frac{d d_\rho^{\pi_{\theta^*}}}{ d d_{\rho}^{\pi_\theta}}\right|^2_{B_b(S)} \left[\int_{S} \lambda_{\min}(G^{\pi_\theta}(s))d_\rho^{\pi_{\theta^*}}(ds)\right]^{-2} \left\|\nabla V_\tau^{\pi_\theta}(\rho)\right\|_2^2, \\
& = \frac{1-\gamma}{2\tau e^{-\hat{R}}}\left|\frac{d d_\rho^{\pi_{\theta^*}}}{ d d_{\rho}^{\pi_\theta}}\right|^2_{B_b(S)}  \lambda^\theta \left\|\nabla V_\tau^{\pi_\theta}(\rho)\right\|_2^2\,. 
\end{align*}
Moreover, note that for any $E \in \mathcal{B}(S)$ and $\pi \in \mathcal{P}(A | S)$, we have
\begin{equation}
\label{eq:kernel_occ_geq_rho}
\begin{split}
d_\rho^\pi(E) & =\int_S d^\pi\left(E | s^{\prime}\right) \rho\left(d s^{\prime}\right)=\int_S(1-\gamma) \sum_{n=0}^{\infty} \gamma^n P_\pi^n\left(E | s^{\prime}\right) \rho\left(d s^{\prime}\right) \\
& \geq \int_S(1-\gamma) P_\pi^0\left(E | s^{\prime}\right) \rho\left(d s^{\prime}\right)=(1-\gamma) \int_S \delta_{s^{\prime}}(E) \rho\left(d s^{\prime}\right)=(1-\gamma) \rho(E)\,. 
\end{split}
\end{equation}
Hence 
\[
V_\tau^{\pi_\theta}(\rho) - V_\tau^{\pi_{\theta^*}}(\rho) 
\leq \frac{1}{2(1-\gamma)\tau e^{-\hat{R}}}\left|\frac{d d_\rho^{\pi_{\theta^*}}}{ d \rho}\right|^2_{B_b(S)}  \lambda^\theta \left\|\nabla V_\tau^{\pi_\theta}(\rho)\right\|_2^2\,.
\]
Thus we have shown that~\eqref{eq:pl_theorem_appendix} holds.

If, additionally,  we know that the feature basis $g:S\times A \to \mathbb R^p$ satisfies $\sum_{i=1}^p g_i(s,a) = 1 $ and for all $s \in S,a \in A$, $g(s,a) \geq 0$ then we return to~\eqref{eq:sandwich_ineq} to proceed.
By Lemma \ref{lem:app:kl-logit_ineq}, we have
\begin{align*}
V_\tau^{\pi_\theta}(\rho) - V_\tau^{\pi_{\theta^*}}(\rho) & \leq \frac{\tau}{1-\gamma}\int \mathrm{KL}(\pi_\theta | \bar{\pi}_{\theta})(s)d_\rho^{\pi_{\theta^*}}(ds)\\
& \leq \frac{\tau }{2(1-\gamma)W}\left\|\theta - \boldsymbol{\theta}(\pi_\theta) - \frac{(\theta - \boldsymbol{\theta}(\pi_\theta))^\top \textbf{1}}{p} \cdot \textbf{1}\right\|_{2}^2,
\end{align*}
where $\hat{R} = \max\left(\frac{2}{(1-\gamma)\tau} \left(1+\gamma\tau R\right), R\right)$.
On the other hand, by~\eqref{eq:grad_for_PL},
\begin{align*}
& \nabla V_\tau^{\pi_\theta}(\rho) \\
& = \frac{\tau}{1-\gamma} \int_{S} G^{\pi_\theta}(s)d_{\rho}^{\pi_{\theta}}(ds)\cdot(\theta -\boldsymbol{\theta}(\pi_\theta))\\
& = \frac{\tau}{1-\gamma} \int_{S} G^{\pi_\theta}(s)d_{\rho}^{\pi_{\theta}}(ds)\cdot\left(\theta -\boldsymbol{\theta}(\pi_\theta)- \frac{(\theta - \boldsymbol{\theta}(\pi_\theta))^\top\textbf{1}}{p}\textbf{1}\right),
\end{align*}
where the last equality is due to Proposition \ref{prp:app:mdp:singularity_of_fim_mdp_simplex}. 
Hence by Corollary \ref{cor:app:mdp:spectrum_of_fim}, we have
{\small
\begin{align*}
& \left\|\nabla V_\tau^{\pi_\theta}(\rho)\right\|_2 \\
& \geq \frac{\tau }{1-\gamma}\int_{S}\lambda_{\min}\left(\int_{A} g(s,a)g^\top(s,a)\pi_\theta(da|s)\right)d^{\pi_{\theta}}_{\rho}(ds)\left\|\theta - \boldsymbol{\theta}(\pi_\theta) - \frac{(\theta - \boldsymbol{\theta}(\pi_\theta))^\top\textbf{1}}{p}\textbf{1}\right\|_2 \\
& = \frac{\tau }{1-\gamma}\int_{S}\lambda_{\min}\left(\int_{A} g(s,a)g^\top(s,a)\pi_\theta(da|s)\right)\frac{d^{\pi_{\theta}}_{\rho}}{d^{\pi_{\theta^*}}_{\rho}}d^{\pi_{\theta^*}}_{\rho}(ds)\left\|\theta - \boldsymbol{\theta}(\pi_\theta) - \frac{(\theta - \boldsymbol{\theta}(\pi_\theta))^\top\textbf{1}}{p}\textbf{1}\right\|_2 \\
& \geq \frac{\tau }{1-\gamma}\left|\frac{\text{d}d^{\pi_{\theta^*}}_{\rho}}{\text{d}d^{\pi_{\theta}}_{\rho}}\right|_{B_b(S)}^{-1}\int_{S}\lambda_{\min}\left(\int_{A} g(s,a)g^\top(s,a)\pi_\theta(da|s)\right)d^{\pi_{\theta^*}}_{\rho}(ds)\left\|\theta - \boldsymbol{\theta}(\pi_\theta) - \frac{(\theta - \boldsymbol{\theta}(\pi_\theta))^\top\textbf{1}}{p}\textbf{1}\right\|_2.
\end{align*}
}
As a result, let
\[
\lambda^{\theta} = \left[\int_{S}\lambda_{\min}\left(\int_{A} g(s,a)g^\top(s,a)\pi_\theta(da|s)\right)d^{\pi_{\theta}}_{\rho}(ds)\right]^{-2}.
\]
We have
\begin{align*}
& V_\tau^{\pi_\theta}(\rho) - V_\tau^{\pi_{\theta^*}}(\rho)\\
& \leq \frac{\tau }{2(1-\gamma) e^{-\hat{R}}}\left\|\theta - \boldsymbol{\theta}(\pi_\theta) - \frac{(\theta - \boldsymbol{\theta}(\pi_\theta))^\top\textbf{1}}{p}\textbf{1}\right\|_{2}^2\\
& \leq \frac{1-\gamma}{2\tau e^{-\hat{R}}}\left|\frac{d d_\rho^{\pi_{\theta^*}}}{ d d_{\rho}^{\pi_\theta}}\right|^2_{B_b(S)} \lambda^\theta \left\|\nabla V_\tau^{\pi_\theta}(\rho)\right\|_2^2 \\
& \leq \frac{1}{2(1-\gamma)\tau e^{-\hat{R}}}\left|\frac{d d_\rho^{\pi_{\theta^*}}}{ d \rho}\right|^2_{B_b(S)} \lambda^\theta \left\|\nabla V_\tau^{\pi_\theta}(\rho)\right\|_2^2 .
\end{align*}
where the last inequality is due to~\eqref{eq:kernel_occ_geq_rho}.This concludes the second situation.
\end{proof}


\subsubsection{Proof of Theorem \ref{thm:main:mdp:Radial_unboundedness_of_relative_entropy}}
\label{sec:pf_Radial_unboundedness_of_relative_entropy}

\aspAffineMaxSet*

\begin{lemma}[Full affine span at a fixed state]\label{rem:app:mdp:affine_span}
For any fixed $s\in S$, the following statements are equivalent:
\begin{enumerate}
\item For any $v \in \mathbb R^p$ such that $\|v\|_2 = 1$ the map $A \ni a \mapsto v^\top g(s,a)$ is not a constant.
\item $g(s,A)$ has full affine dimension i.e. $\operatorname{span}\{g(s,a)-g(s,a'):a,a'\in A\}=\mathbb R^p$.
\end{enumerate}
\end{lemma}

\begin{proof}
To see this, let us first show that 1. implies 2.
We proceed by contradiction.
If the span in 2. was a proper subspace $V_s \subset \mathbb R^p$,
then $V_s^\perp$ is not empty. 
Hence we can choose some $m\in V_s^\perp\setminus\{\mathbf 0\}$ such that
\[
m^\top(g(s,a)-g(s,a'))=0,\qquad \forall a,a'\in A.
\]
But then $a\mapsto (m/\|m\|_2)^\top g(s,a) = (m/\|m\|_2)^\top g(s,a')$ which is a constant thus contradicting 1.

Now we show 2. implies 1.
Again we proceed by contradiction.
From 2. we know that $u\in V_s^\perp=\{\mathbf 0\}$.
If $u^\top g(s,\cdot)$ is constant for some
$u\in \mathbb R^p$ such that $\|u\|_2 = 1$ then $u^\top(g(s,a)-g(s,a'))=0$ for all
$a,a'\in A$.
So $u\in V_s^\perp=\{\mathbf 0\}$, a contradiction. 
\end{proof}


\KLUBAFFINE*

\begin{proof}            
Consider $\theta = r u$ with $r = \|\theta\|_2$ and $u \in \mathbb{S}^{p-1}$. 
We analyze along $ \theta = r u $ with $ r \to +\infty $:
Define for fixed $s \in S$,
\[
\phi^s_u(a) := u^\top g(s,a),
\qquad
\phi_u^{\max}(s) := \max_{a \in A} \phi^s_u(a).
\]

Then
\[
\theta^\top g(s,a) = r \phi^s_u(a),
\qquad
Z(\theta) = \int_{A} e^{r \phi^s_u(a)} \mu(da).
\]

By assumption,
\[
\mu\bigl(\operatorname{Argmax}(\phi^s_u)\bigr) = 0.
\]

For any fixed margin $\varepsilon > 0$, define the $\varepsilon$-optimal super-level set:
\[
A_\varepsilon := \{ a \in A : \phi^s_u(a) > \phi_u^{\max}(s) - \varepsilon \}.
\]

First, we establish that the Gibbs measure strictly concentrates on $A_\varepsilon$. 
The probability of sampling outside this set under $\pi_\theta$ is bounded by evaluating the worst-case numerator and restricting the denominator's domain to $A_{\varepsilon/2}$:
\[
\pi_{ru}(A_\varepsilon^c | s) = \frac{\int_{A_\varepsilon^c} e^{r \phi^s_u(a)} \mu(da)}{\int_A e^{r \phi^s_u(a)} \mu(da)} \le \frac{e^{r(\phi_u^{\max}(s) - \varepsilon)}}{\int_{A_{\varepsilon/2}} e^{r \phi^s_u(a)} \mu(da)} \le \frac{e^{r(\phi_u^{\max}(s) - \varepsilon)}}{e^{r(\phi_u^{\max}(s) - \varepsilon/2)} \mu(A_{\varepsilon/2})} = \frac{e^{-r\varepsilon/2}}{\mu(A_{\varepsilon/2})}.
\]

$\mu(A_{\varepsilon/2}) > 0$ is a strictly positive constant independent of $r$
And we have $e^{-r\varepsilon/2} \to 0$ as $r \to \infty$. Thus, $\pi_{ru}(A_\varepsilon^c | s) \to 0$, which implies $\pi_{ru}(A_\varepsilon | s) \to 1$.

Next, we apply partition property of the KL divergence (Lemma 1.4.3 in~\cite{dupuis2011weak}). 
By splitting the action space into the binary partition $\{A_\varepsilon, A_\varepsilon^c\}$, the KL divergence is bounded below by the divergence between the Bernoulli distributions induced by this partition:
\[
\mathrm{KL}(\pi_{ru}(\cdot |s) | \mu) \ge \pi_{ru}(A_\varepsilon | s) \log \frac{\pi_{ru}(A_\varepsilon | s)}{\mu(A_\varepsilon)} + \pi_{ru}(A_\varepsilon^c | s) \log \frac{\pi_{ru}(A_\varepsilon^c | s)}{\mu(A_\varepsilon^c)}.
\]

Taking the limit inferior as $r \to \infty$, and substituting $\pi_{ru}(A_\varepsilon | s) \to 1$ and $\pi_{ru}(A_\varepsilon^c | s) \to 0$, we obtain:
\[
\liminf_{r \to \infty} \mathrm{KL}(\pi_{ru}(\cdot |s) | \mu) \ge 1 \cdot \log \frac{1}{\mu(A_\varepsilon)} + 0 = -\log \mu(A_\varepsilon).
\]

This lower bound holds for any $\varepsilon > 0$. By the continuity of $g(s,\cdot)$, the intersection of these sets as $\varepsilon \downarrow 0$ is exactly the maximizer set $\operatorname{Argmax}(\phi^s_u)$ i.e.
\[
\bigcap_{\varepsilon > 0} A_\varepsilon
=
\operatorname{Argmax}(\phi^s_u).
\]

Moreover, the family $\{A_\varepsilon\}_{\varepsilon > 0}$ is decreasing as
$\varepsilon \downarrow 0$.
By continuity from above of the probability measure $\mu$,
\[
\lim_{\varepsilon \downarrow 0} \mu(A_\varepsilon)
=
\mu\!\left( \bigcap_{\varepsilon > 0} A_\varepsilon \right)
=
\mu(\operatorname{Argmax}(\phi^s_u)).
\]

By Assumption \ref{asp:app:mdp:max_set_affine}, $\mu(\operatorname{Argmax}(\phi^s_u)) = 0$. Therefore, taking the limit as $\varepsilon \downarrow 0$ yields:
\[
\lim_{\varepsilon \downarrow 0} (-\log \mu(A_\varepsilon)) = +\infty.
\]

Consequently, since $u \in \mathbb{S}^{p-1}$ is arbitrary and $\mathbb{S}^{p-1}$ is compact, $\lim_{\|\theta\|_2\to\infty}
\mathrm{KL}\bigl(\pi_\theta(\cdot\mid s)\,|\,\mu\bigr)=\infty$, which completes the proof.
\end{proof}


\subsubsection{Proof of Lemma \ref{lem:main:mdp:exist_sol_gf_affine}}
\label{sec:pf_exist_sol_gf_affine}
\ExistSolAffine*

\begin{proof}
Since the gradient flow satisfying local Lipschitz condition by Proposition \ref{prp:app:mdp:local_lipchitiz_conitinuity_mdp}, with Lyapunov function $V_\tau^{\pi_{\cdot}}(\rho): \mathbb{R}^p \rightarrow \mathbb{R}$ that has radial unboundedness by Theorem \ref{thm:main:mdp:Radial_unboundedness_of_relative_entropy}, and
\[
\frac{d}{dt}V_\tau^{\pi_{\theta_t}}(\rho)=-\left\|\nabla V_\tau^{\pi_{\theta_t}}(\rho)\right\|_2^2 \leq 0,
\]
we know the existence of solutions to the gradient flow for any $ t \geq 0$ and there exists a constant $C_\tau$ such that $\sup_t\|\theta_t\| \leq C_\tau $ by classical arguments for constructing ODE solutions with Lyapunov functions.

For all $s \in S, a \in A$, any $t \geq 0$
\begin{align*}
&\left|\log\frac{d \pi_{\theta_t}}{d \mu}(a|s)\right|\\
& = \left| \langle \theta_t, g(s,a)\rangle - \log Z_{\pi_{\theta_t}}(s)\right| \\
& \leq \sup_{s,a}\|g(s,a)\|_2\|\theta_t\|_2 + \left|\log \int_{A} e^{\langle \theta_t,g(s,a)\rangle}\mu(da)\right|\\
& \leq 2C_\tau < \infty,
\end{align*}
\end{proof}

\subsubsection{Proof of Theorem \ref{thm:main:lin_cov_affine}}
\label{sec:pf_lin_cov_affine}
\aspInitAffine* 
\LinCovAffine*  

\begin{proof}
From Lemma \ref{lem:main:mdp:exist_sol_gf_affine}, there exists $R >0$ which is sufficiently large, such that for $t \geq 0$, 
\[
\theta^*, \theta_t \in \Theta_{R} = \{\theta: \pi_\theta \in \Pi_\mu \ \text{and}\  |\log \frac{d \pi_\theta}{d\mu}|_{B_b(S\times A)} \leq R \},
\] 
which means that $\min_{s,a}\frac{d \pi_{\theta_t}}{d \mu}(a|s) \geq e^{-R}$.
From Lemma \ref{lemapp:mdp:lower_bound_log_density}, we know that there exists 
\[
\lambda = e^{-R-2\|\theta_0\|_2}\int_{S}   \lambda_{\min} \left(G^{\pi_{\theta_0}}(s)\right) d_\rho^{\pi^*}(ds) >0
\] 
such that 
\begin{equation}\label{eq:positive_lambda_affine}
\int_{S}\lambda_{\min} (G^{\pi_{\theta_t}}(s))d_\rho^{\pi^*}(ds) \geq \lambda >0.
\end{equation}

From Theorem \ref{thm:main: non-uniform_PL_ineq}, we know
\begin{align*}
& \sup_{t \geq 0} C_R(\theta_t) \nonumber\\
& = \sup_{t \geq 0} \frac{1}{2\tau(1-\gamma) \exp\left({-\max\left(\frac{2}{(1-\gamma)\tau} \left(1+\gamma\tau R\right), R\right)}\right)}\left|\frac{d d_\rho^{\pi^*}}{ d \rho}\right|^2_{B_b(S)} \bigg[\int_S \lambda_{\min}(G^{\pi_{\theta_t}}(s))d_\rho^{\pi^*}(ds)\bigg]^{-2}\\
& \leq \frac{1}{2\tau \lambda^2 (1-\gamma) \exp\left({-\max\left(\frac{2}{(1-\gamma)\tau} \left(1+\gamma\tau R\right), R\right)}\right)}\left|\frac{d d_\rho^{\pi^*}}{ d \rho}\right|^2_{B_b(S)} := C_{\theta_0}.
\end{align*}

Then from Proposition \ref{prp:app:mdp:Gradient_of_the_Objective}, and Theorem \ref{thm:main: non-uniform_PL_ineq} again,
\begin{align*}
\frac{d}{dt} V_\tau^{\pi_{\theta_t}}(\rho) &= - \left\|\nabla V_\tau^{\pi_{\theta_t}}(\rho)\right\|^2_2\\
& \leq -C_{\theta_0}^{-1} \left(V_\tau^{\pi_{\theta_t}}(\rho) - V_\tau^{\pi^*}  (\rho) \right)
\end{align*}
then using Grönwall's inequality concludes the proof.
\end{proof}

\subsubsection{Proof of Theorem \ref{thm:main:mdp:Radial_unboundedness_of_relative_entropy_in_subspace_simplex}}
\label{sec:pf_Radial_unboundedness_of_relative_entropy_in_subspace_simplex}

\aspSimplexMaxSet*

Recall that for any $v\in \mathbb R^p$ we will write $v_\perp := v - p^{-1}\langle v, \mathbf 1\rangle \mathbf 1$.

\KLUBSIMPLEX*
\begin{proof}
Consider $\theta_\perp = r u$ with $r = \|\theta\|_2$ and $u \perp \mathbf 1, u \in \mathbb{S}^{p-1}$. 
We analyze along $ \theta_\perp = r u $ with $ r \to +\infty $:
Define for fixed $s \in S$,
\[
\phi^s_u(a) := u^\top g(s,a),
\qquad
\phi_u^{\max}(s) := \max_{a \in A} \phi^s_u(a).
\]

Then
\[
\theta^\top g(s,a) = r \phi^s_u(a),
\qquad
Z(\theta) = \int_{A} e^{r \phi^s_u(a)} \mu(da).
\]

By assumption,
\[
\mu\bigl(\operatorname{Argmax}(\phi^s_u)\bigr) = 0.
\]

For any fixed margin $\varepsilon > 0$, define the $\varepsilon$-optimal super-level set:
\[
A_\varepsilon := \{ a \in A : \phi^s_u(a) > \phi_u^{\max}(s) - \varepsilon \}.
\]

First, we establish that the Gibbs measure strictly concentrates on $A_\varepsilon$. 
The probability of sampling outside this set under $\pi_\theta$ is bounded by evaluating the worst-case numerator and restricting the denominator's domain to $A_{\varepsilon/2}$:
\[
\pi_{ru}(A_\varepsilon^c | s) = \frac{\int_{A_\varepsilon^c} e^{r \phi^s_u(a)} \mu(da)}{\int_A e^{r \phi^s_u(a)} \mu(da)} \le \frac{e^{r(\phi_u^{\max}(s) - \varepsilon)}}{\int_{A_{\varepsilon/2}} e^{r \phi^s_u(a)} \mu(da)} \le \frac{e^{r(\phi_u^{\max}(s) - \varepsilon)}}{e^{r(\phi_u^{\max}(s) - \varepsilon/2)} \mu(A_{\varepsilon/2})} = \frac{e^{-r\varepsilon/2}}{\mu(A_{\varepsilon/2})}.
\]

$\mu(A_{\varepsilon/2}) > 0$ is a strictly positive constant independent of $r$
And we have $e^{-r\varepsilon/2} \to 0$ as $r \to \infty$. Thus, $\pi_{ru}(A_\varepsilon^c | s) \to 0$, which implies $\pi_{ru}(A_\varepsilon | s) \to 1$.

Next, we apply partition property of the KL divergence (Lemma 1.4.3 in~\cite{dupuis2011weak}). 
By splitting the action space into the binary partition $\{A_\varepsilon, A_\varepsilon^c\}$, the KL divergence is bounded below by the divergence between the Bernoulli distributions induced by this partition:
\[
\mathrm{KL}(\pi_{ru}(\cdot |s) | \mu) \ge \pi_{ru}(A_\varepsilon | s) \log \frac{\pi_{ru}(A_\varepsilon | s)}{\mu(A_\varepsilon)} + \pi_{ru}(A_\varepsilon^c | s) \log \frac{\pi_{ru}(A_\varepsilon^c | s)}{\mu(A_\varepsilon^c)}.
\]

Taking the limit inferior as $r \to \infty$, and substituting $\pi_{ru}(A_\varepsilon | s) \to 1$ and $\pi_{ru}(A_\varepsilon^c | s) \to 0$, we obtain:
\[
\liminf_{r \to \infty} \mathrm{KL}(\pi_{ru}(\cdot |s) | \mu) \ge 1 \cdot \log \frac{1}{\mu(A_\varepsilon)} + 0 = -\log \mu(A_\varepsilon).
\]

This lower bound holds for any $\varepsilon > 0$. By the continuity of $g(s,\cdot)$, the intersection of these sets as $\varepsilon \downarrow 0$ is exactly the maximizer set $\operatorname{Argmax}(\phi^s_u)$ i.e.
\[
\bigcap_{\varepsilon > 0} A_\varepsilon
=
\operatorname{Argmax}(\phi^s_u).
\]

Moreover, the family $\{A_\varepsilon\}_{\varepsilon > 0}$ is decreasing as
$\varepsilon \downarrow 0$.
By continuity from above of the probability measure $\mu$,
\[
\lim_{\varepsilon \downarrow 0} \mu(A_\varepsilon)
=
\mu\!\left( \bigcap_{\varepsilon > 0} A_\varepsilon \right)
=
\mu(\operatorname{Argmax}(\phi^s_u)).
\]

By Assumption \ref{asp:app:mdp:max_set_affine}, $\mu(\operatorname{Argmax}(\phi^s_u)) = 0$. Therefore, taking the limit as $\varepsilon \downarrow 0$ yields:
\[
\lim_{\varepsilon \downarrow 0} (-\log \mu(A_\varepsilon)) = +\infty.
\]

Consequently, since $u \in \mathbb{S}^{p-1} \cap \{v \in \mathbb{R}^p: v \perp \mathbf 1\}$ is arbitrary and $\mathbb{S}^{p-1} \cap \{v \in \mathbb{R}^p: v \perp \mathbf 1\}$ is compact, $\lim_{\theta\in \mathbb R^p:\|\theta_\perp\|_2\to\infty}
\mathrm{KL}\bigl(\pi_\theta(\cdot\mid s)\,|\,\mu\bigr)=\infty$, which completes the proof.
\end{proof}

\subsubsection{Proof of Lemma \ref{lem:main:mdp:exist_sol_gf_simplex}}
\label{sec:pf_exist_sol_gf_simplex}

\ExistSolSimplex*

\begin{proof}

By Proposition~\ref{prp:app:mdp:local_lipchitiz_conitinuity_mdp}, we know the solution to~\eqref{eq:intro:continuous_GF} exists for $t \in [0, T_{\max})$  for some $T_{\max} < \infty$, denoted by $\{\theta_t\}_{t \in [0, T_{\max})}$.
Firstly, we show that the solution to the gradient flow would not explode in finite time $t \in [0, T_{\max})$. Note that for all $s \in S$ and $a \in A$,
\begin{align}
&\left|\log\frac{d \pi_{\theta_t}}{d \mu}(a|s)\right| \nonumber\\
& = \left| \langle \theta_t, g(s,a)\rangle - \log Z_{\pi_{\theta_t}} \right| \nonumber\\
& \leq \sup_{s,a}\|g(s,a)\|_2\|\theta_t\|_2 + \left|\log \int_A e^{\langle \theta_t,g(s,a)\rangle}\mu(da)\right| \nonumber\\
& \leq 2  \|\theta_t\|_2 \label{eq:log_density_bound_*}.
\end{align}

Calculation in Lemma \ref{lem:bounds} shows that
\begin{align*}
|V^{\pi_\theta}_{\tau}(s)|&=\frac{1}{1-\gamma}\left|\int_{S}\int_A \left(c(s',a)+\tau \log \frac{d\pi_{\theta}}{d\mu}(a|s)\right)\pi_{\theta}(da|s')d^{\pi_\theta}(ds'|s)\right|\\
&\le \frac{1}{1-\gamma} \left(1+\tau \left|\log \frac{d\pi_{\theta}}{d\mu}\right|_{B_b(S \times A)}\right)\,.
\end{align*}
and
\begin{equation}\label{eq:Qpi_bound_*}
|Q^{\pi_{\theta}}_{\tau}(s,a)| \le \frac{1}{1-\gamma} \left(1+\gamma\tau  \left|\log \frac{d\pi_{\theta}}{d\mu}\right|_{B_b(S \times A)}\right)\,.
\end{equation}
By Proposition \ref{prp:app:mdp:Gradient_of_the_Objective}, we have

\begin{align*}
\frac{1}{2}\frac{d}{d t}\|\theta_t\|_2^2
& = \frac{1}{1-\gamma}\left \langle -\tau G^{\pi_{\theta_t}}(\rho)  \theta_t -  \int_{S} \int_{A} Q^{\pi_{\theta_t}}_{\tau}(s,a)\nabla\log\frac{d \pi_{\theta_t}}{d \mu} (a|s)\pi_{\theta_t}(da|s)d_\rho^{\pi_t}(ds), \theta_t\right \rangle \\
& \leq \frac{1}{1-\gamma}\left \langle  -  \int_{S} \int_{A} Q^{\pi_{\theta_t}}_{\tau}(s,a)\nabla\log\frac{d \pi_{\theta_t}}{d \mu} (a|s)\pi_{\theta_t}(da|s)d_\rho^{\pi_t}(ds), \theta_t\right \rangle \\
& \leq \frac{1}{2(1-\gamma)} \|\theta_t\|_2^2 + \frac{1}{2(1-\gamma)}\left(\frac{1}{1-\gamma} \left(1+\gamma\tau \left|\log \frac{d\pi_{\theta_t}}{d\mu}\right|_{B_b(S \times A)}\right) \cdot 2\right)^2\\
& \leq \frac{1}{2(1-\gamma)} \|\theta_t\|_2^2 + \frac{1}{2(1-\gamma)}\left(\frac{2}{1-\gamma} \left(1+\gamma\tau 2\|\theta_t\|_2 \right)\right)^2.
\end{align*}
where in the first inequality we use that  $G^{\pi_{\theta_t}}(\rho)$ is positive semi-definite, in the second inequality we use Young's Inequality, bounds~\eqref{eq:log_density_bound_*} and Lemma~\ref{lem:app:mdp:bounded_grad_of_log_policy}, in the last inequality we use~\eqref{eq:Qpi_bound_*}. From the previous estimates, define
\[
y(t) := \|\theta_t\|^2_2 .
\]
Then we obtain the differential inequality
\[
\frac{d}{dt} y(t) \le C_1 y(t) + C_2 ,
\]
for constants $C_1, C_2 > 0$. By Grönwall’s inequality,
\[
y(t) \le e^{C_1 t} y(0) + \frac{C_2}{C_1}\big(e^{C_1 t} - 1\big).
\]
Equivalently,
\[
\|\theta_t\|_2 \le \sqrt{ e^{C_1 t}\|\theta_0\|^2_2 + \frac{C_2}{C_1}\big(e^{C_1 t} - 1\big) }.
\] 
As a result, there is no finite time blow up of $\theta_t$. 

Under Assumption \ref{asp:app:mdp:max_set_simplex}, by Proposition \ref{prp:app:mdp:Gradient_of_the_Objective} and Lemma \ref{lem:app:mdp:grad_of_log_policy}, we know that
\begin{align*}
\frac{d}{d t}\left(\mathbf{1}^{\top} \theta_t\right) &= \frac{1}{1-\gamma}\int\int \left(Q^{\pi_\theta}_\tau(s,a) + \tau \log \frac{d \pi_{\theta}}{d \mu}(a|s)\right) \mathbf{1}^\top\nabla \log \frac{d\pi_\theta}{d \mu}(a| s)\pi_{\theta}(da|s)d_\rho^{\pi_{\theta}}(ds)\\
& = 0,
\end{align*}
implying that $\mathbf{1}^\top \theta_t = \mathbf{1}^\top \theta_0$ for any $t \geq 0$.

Consider $\theta_t = (\theta_{t})_{\perp} + \frac{\theta_t^\top \mathbf{1}}{p}\mathbf{1}$, where $\theta_{t,\perp} \perp \mathbf{1}$.
By Theorem \ref{thm:main:mdp:Radial_unboundedness_of_relative_entropy_in_subspace_simplex}, and the fact 
\[
\frac{d}{dt}V_\tau^{\pi_{\theta_t}}(\rho)=-\left\|\nabla V_\tau^{\pi_{\theta_t}}(\rho)\right\|_2^2 \leq 0,
\]
we know $ S=\Omega_c \cap\left\{\theta | \mathbf{1}^T \theta=\mathbf{1}^T \theta_0\right\}$, where $\Omega_c = \left\{\theta | V_\tau^{\pi_\theta}(\rho) \leq V_\tau^{\pi_{\theta_0}}(\rho)\right\}$, is compact positive invariant. 
Hence classical arguments for constructing ODE solutions with Lyapunov functions, we know the existence of solutions to the gradient flow for any $ t \geq 0$ and there exists a constant $C_\tau$ such that $\sup_t\|\theta_t\| \leq C_\tau $ .

For all $s \in S, a \in A$, any $t \geq 0$,
\begin{align*}
&\left|\log\frac{d \pi_{\theta_t}}{d \mu}(a|s)\right|\\
& = \left| \langle \theta_t, g(s,a)\rangle - \log Z_{\pi_{\theta_t}}(s)\right| \\
& \leq \sup_{s,a}\|g(s,a)\|_2\|\theta_t\|_2 + \left|\log \int_{A} e^{\langle \theta_t,g(s,a)\rangle}\mu(da)\right|\\
& \leq 2C_\tau < \infty,
\end{align*}
\end{proof}

\subsubsection{Proof of Theorem  \ref{thm:main:lin_cov_simplex}}
\label{sec:pf_lin_cov_simplex}

\aspInitSimplex*
\LinCovSimplex*

\begin{proof}
From Lemma \ref{lem:main:mdp:exist_sol_gf_simplex}, there exists $R >0$ which is sufficiently large, such that for $t \geq 0$, 
\[
\theta^*, \theta_t \in \Theta_{R} = \{\theta: \pi_\theta \in \Pi_\mu \ \text{and}\  |\log \frac{d \pi_\theta}{d\mu}|_{B_b(S\times A)} \leq R \}.
\]
which means that $\min_{s,a}\frac{d \pi_{\theta_t}}{d \mu}(a|s) \geq e^{-R}$.
From Lemma \ref{lem:app:mdp:lower_bound_log_density_simplex}, we know that there exists 
\begin{equation}\label{eq:positive_lambda_simplex}
\lambda = e^{-R-2\|\theta_0\|_2}  \int_{S} \lambda_{\min}\left(\int_{A} g(s,a)g^\top(s,a)\pi_{\theta_0} (da | s)\right) d_\rho^{\pi^*}(ds) >0
\end{equation}
such that 
\[
\int_{S} \lambda_{\min} (\int_{A} g(s,a)g^\top(s,a)\pi_{\theta_t} (da | s))d_\rho^{\pi^*}(ds) \geq \lambda >0.
\]

From Theorem \ref{thm:main: non-uniform_PL_ineq}, we know
\begin{align*}
& \sup_{t \geq 0} C_R(\theta_t) \nonumber\\
& = \sup_{t \geq 0} \frac{1}{2\tau (1-\gamma) \exp\left({-\max\left(\frac{2}{(1-\gamma)\tau} \left(1+\gamma\tau R\right), R\right)}\right)}\left|\frac{d d_\rho^{\pi^*}}{ d \rho}\right|^2_{B_b(S)} \\
& \quad \quad \quad \cdot\left[\int_{S}\lambda_{\min}\left(\int_{A} g(s,a)g^\top(s,a)\pi_\theta(da|s)\right)d^{\pi_{\theta}}_{\rho}(ds)\right]^{-2}\\
& \leq \frac{1}{2\tau \lambda^2 (1-\gamma) \exp\left({-\max\left(\frac{2}{(1-\gamma)\tau} \left(1+\gamma\tau R\right), R\right)}\right)}\left|\frac{d d_\rho^{\pi^*}}{ d \rho}\right|^2_{B_b(S)} := C_{\theta_0}.
\end{align*}

Then from Proposition \ref{prp:app:mdp:Gradient_of_the_Objective}, and Theorem \ref{thm:main: non-uniform_PL_ineq} again,
\begin{align*}
\frac{d}{dt} V_\tau^{\pi_{\theta_t}}(\rho) &= - \left\|\nabla V_\tau^{\pi_{\theta_t}}(\rho)\right\|^2_2\\
& \leq -C_{\theta_0}^{-1} \left(V_\tau^{\pi_{\theta_t}}(\rho) - V_\tau^{\pi^*}  (\rho) \right)
\end{align*}
then using Grönwall's inequality concludes the proof.
\end{proof}

\subsection{More discussions about Example~\ref{ex:hat_func_short}}
\label{sec:ex}
\begin{example}[Hat function features]\label{ex:hat_func}
Let $x_0<x_1<\cdots <x_N$, and consider the standard one-dimensional finite-element hat functions
\[
g_i(a)=
\begin{cases}
\dfrac{a-x_{i-1}}{x_i-x_{i-1}}, & a\in [x_{i-1},x_i],\\[0.8em]
\dfrac{x_{i+1}-a}{x_{i+1}-x_i}, & a\in [x_i,x_{i+1}],\\[0.8em]
0, & \text{otherwise},
\end{cases}
\qquad 1\leq i\leq N-1,
\]
together with the boundary functions
\[
g_0(a)=
\begin{cases}
\dfrac{x_1-a}{x_1-x_0}, & a\in [x_0,x_1],\\[0.8em]
0, & \text{otherwise},
\end{cases}
\qquad
g_N(a)=
\begin{cases}
\dfrac{a-x_{N-1}}{x_N-x_{N-1}}, & a\in [x_{N-1},x_N],\\[0.8em]
0, & \text{otherwise}.
\end{cases}
\]
Set $g(a):=(g_0(a),g_1(a),\dots,g_N(a))^\top$. We now take $N=3$ and choose the concrete grid $x_0=0,\ x_1=\frac13,\ x_2=\frac23,\ x_3=1$.

Let $\theta=(-1,1,1,-1)^\top$. Then $\theta^\top \mathbf 1=-1+1+1-1=0$. Define $f_\theta(a):=\theta^\top g(a)$. Therefore
\[
f_\theta(a)=
\begin{cases}
6a-1, & a\in \left[0,\frac13\right],\\[0.8em]
1, & a\in \left[\frac13,\frac23\right],\\[0.8em]
5-6a, & a\in \left[\frac23,1\right].
\end{cases}
\]
Thus $f_\theta(a)\leq 1$ for every $a\in[0,1]$, and the maximizer set is $\operatorname*{arg\,max}_{a\in[0,1]} f_\theta(a)=\left[\frac13,\frac23\right]$, which has positive Lebesgue measure $\mu([\frac13,\frac23])=\frac13$. Consequently, Assumption~\ref{asp:app:mdp:max_set_simplex} is not satisfied in this example.

Let $\mu$ be Lebesgue measure restricted to $[0,1]$. For $\beta>0$, define the probability measure $\pi_\beta$ by $p_\beta(a):=\frac{d\pi_\beta}{d\mu}(a)=e^{\beta f_\theta(a)}/Z_\beta$, where $Z_\beta:=\int_0^1 e^{\beta f_\theta(a)}\,da$. Since $p_\beta$ integrates to one, the relative entropy of $\pi_\beta$ with respect to $\mu$ is $\operatorname{KL}(\pi_\beta| \mu):=\int_0^1 p_\beta(a)\log p_\beta(a)\,da$. Moreover, since $\log p_\beta(a)=\beta f_\theta(a)-\log Z_\beta$, we have
\[
\operatorname{KL}(\pi_\beta| \mu)
=
\beta\,\mathbb E_{\pi_\beta}[f_\theta(a)]-\log Z_\beta .
\]

In the present grid, $L_1:=x_1-x_0=\frac13$, $L:=x_2-x_1=\frac13$, $L_2:=x_3-x_2=\frac13$.

We first compute $Z_\beta$. On $[0,\frac13]$, set $y=f_\theta(a)=6a-1$. Then $\int_0^{1/3} e^{\beta f_\theta(a)}\,da=\frac16\int_{-1}^1 e^{\beta y}\,dy$. On the middle interval $[\frac13,\frac23]$, one has $f_\theta(a)=1$, so $\int_{1/3}^{2/3} e^{\beta f_\theta(a)}\,da=\frac13 e^\beta$. On $[\frac23,1]$, set $y=f_\theta(a)=5-6a$. Then $dy=-6\,da$, $\int_{2/3}^{1} e^{\beta f_\theta(a)}\,da=\frac16\int_{-1}^1 e^{\beta y}\,dy$. Combining the three pieces gives
\[
Z_\beta
=
\frac13\int_{-1}^1 e^{\beta y}\,dy+\frac13 e^\beta .
\]
Since $\int_{-1}^1 e^{\beta y}\,dy=(e^\beta-e^{-\beta})/\beta$, we obtain
\[
Z_\beta
=
\frac13\frac{e^\beta-e^{-\beta}}{\beta}+\frac13 e^\beta
=
\frac{e^\beta}{3}\left[1+\frac{1-e^{-2\beta}}{\beta}\right].
\]
Therefore $\log Z_\beta=\beta-\log 3+\log(1+(1-e^{-2\beta})/\beta)$. Since $e^{-2\beta}$ is exponentially small as $\beta\to\infty$, and since $\log(1+u)=u-\frac{u^2}{2}+\mathcal O(u^3)$ as $u\to 0$, we get
\[
\log Z_\beta
=
\beta-\log 3+\frac1\beta-\frac{1}{2\beta^2}
+
\mathcal O\!\left(\frac1{\beta^3}\right).
\]

Next define $N_\beta:=\int_0^1 f_\theta(a)e^{\beta f_\theta(a)}\,da$. Then $\mathbb E_{\pi_\beta}[f_\theta(a)]=N_\beta/Z_\beta$. We compute $N_\beta$ using the same changes of variables. On the two exterior intervals, the variable $y=f_\theta(a)$ runs linearly from $-1$ to $1$, while on the middle interval $f_\theta(a)=1$. Hence
\[
N_\beta
=
\frac13\int_{-1}^1 y e^{\beta y}\,dy+\frac13 e^\beta .
\]
Moreover,
\[
\int_{-1}^1 y e^{\beta y}\,dy
=
\left[
\frac{y e^{\beta y}}{\beta}
-
\frac{e^{\beta y}}{\beta^2}
\right]_{-1}^{1}
=
\frac{e^\beta}{\beta}
-
\frac{e^\beta}{\beta^2}
+
\frac{e^{-\beta}}{\beta}
+
\frac{e^{-\beta}}{\beta^2}.
\]
Thus
\[
N_\beta
=
\frac{e^\beta}{3}
\left[
1+\frac1\beta-\frac1{\beta^2}
+
e^{-2\beta}
\left(
\frac1\beta+\frac1{\beta^2}
\right)
\right].
\]

Dividing the exact formula for $N_\beta$ by the exact formula for $Z_\beta$, we find
\[
\mathbb E_{\pi_\beta}[f_\theta(a)]
=
\frac{
1+\dfrac1\beta-\dfrac1{\beta^2}
+
e^{-2\beta}\left(\dfrac1\beta+\dfrac1{\beta^2}\right)
}{
1+\dfrac{1-e^{-2\beta}}{\beta}
}.
\]
Since $e^{-2\beta}$ is exponentially small, it may be absorbed into any algebraic remainder. Therefore
\[
\mathbb E_{\pi_\beta}[f_\theta(a)]
=
\frac{1+\dfrac1\beta-\dfrac1{\beta^2}}{1+\dfrac1\beta}
+
\mathcal O(e^{-2\beta}).
\]
The rational term equals $1-\frac{1/\beta^2}{1+1/\beta}$. Using $1/(1+1/\beta)=1-1/\beta+\mathcal O(1/\beta^2)$, we obtain
\[
\mathbb E_{\pi_\beta}[f_\theta(a)]
=
1-\frac1{\beta^2}+\frac1{\beta^3}
+
\mathcal O\!\left(\frac1{\beta^4}\right),
\]
and consequently
\[
\beta\,\mathbb E_{\pi_\beta}[f_\theta(a)]
=
\beta-\frac1\beta+\frac1{\beta^2}
+
\mathcal O\!\left(\frac1{\beta^3}\right).
\]

Substituting the expansions of $\beta\,\mathbb E_{\pi_\beta}[f_\theta(a)]$ and $\log Z_\beta$ into the entropy identity gives
\[
\begin{aligned}
\operatorname{KL}(\pi_\beta| \mu)
&=
\beta\,\mathbb E_{\pi_\beta}[f_\theta(a)]
-
\log Z_\beta \\[0.4em]
&=
\log 3
-
\frac{2}{\beta}
+
\frac{3}{2\beta^2}
+
\mathcal O\!\left(\frac1{\beta^3}\right).
\end{aligned}
\]
In particular,
\[
\operatorname{KL}(\pi_\beta| \mu)
=
\log 3
-
\frac{2}{\beta}
+
\mathcal O\!\left(\frac1{\beta^2}\right).
\]
Therefore, $\lim_{\beta\to\infty}\operatorname{KL}(\pi_\beta| \mu)=\log 3$.
Hence, in this concrete example, the relative entropy remains bounded as $\beta\to\infty$. Therefore, without Assumption~\ref{asp:app:mdp:max_set_simplex}, one cannot in general expect radial unboundedness of the KL divergence.
\end{example}

\subsection{Additional useful results}
\label{sec:extra results}

We introduce the following dynamic programming principle (see Theorem B.1 in~\cite{kerimkulov2025fisher}).
\begin{theorem}[Dynamic programming principle]\label{thm:app:dpp}
Let $\tau>0$. The optimal value function $V_\tau^*$ is the unique bounded solution of the following Bellman equation:
\[
V_\tau^*(s)=\inf _{m \in \mathcal{P}(A)} \int_A\left(c(s, a)+\tau \log \frac{d m}{\mathrm{~d} \mu}(a)+\gamma \int_S V_\tau^*\left(s^{\prime}\right) P\left(d s^{\prime} | s, a\right)\right) m(d a), \quad \forall s \in S,
\]
Consequently, for all $s \in S$,
\[
V_\tau^*(s)=-\tau \log \int_A \exp \left(-\frac{1}{\tau} Q_\tau^*(s, a)\right) \mu(d a)
\]
where $Q^* \in B_b(S \times A)$ is defined by
\[
Q_\tau^*(s, a)=c(s, a)+\gamma \int_S V_\tau^*\left(s^{\prime}\right) P\left(d s^{\prime} | s, a\right), \quad \forall(s, a) \in S \times A
\]
Moreover, there is an optimal policy $\pi_\tau^* \in \mathcal{P}_\mu(a|s)$ given by
\[
\pi_\tau^*(d a|s)=\exp \left(-\left(Q_\tau^*(s, a)-V_\tau^*(s)\right) / \tau\right) \mu(d a), \quad \forall s \in S
\]
\end{theorem}


\begin{definition}\label{def:app:flat_derivative}
A functional $F: \mathcal{C} \mapsto \mathbb{R}^d$ is said to admit a linear derivative if there is a continuous map $\frac{\delta F}{\delta m}: \mathcal{C} \times \mathbb{R}^p \mapsto \mathbb{R}^d$, such that for all $m, m^{\prime} \in \mathcal{C}$, it holds that $\int \left\|\frac{\delta F}{\delta m}( m)(a)\right\|_2 m^{\prime}(a) d(a) < \infty$, and
\begin{equation}\label{eq: flat derivative01 }
F\left(m^{\prime}\right)-F(m)=\int_0^1 \int \frac{\delta F}{\delta m}\left(m+\lambda\left(m^{\prime}-m\right)\right)(a)  \cdot \left(m^{\prime}(a)-m(a)\right)da d \lambda .
\end{equation}
\end{definition}

\begin{proposition}[Strong convexity of negative entropy]
\label{prp:app:Strong_convexity_of_negative_entropy}
Let $W>0$ and define
\[
\mathcal{P}_W^{ A}
:=\left\{
\pi\in\mathcal P( A):
\pi\ll\mu,\ 
\left|\frac{d\pi}{d\mu}\right|_{ B_b(A)}
\le W^{-1}
\right\}.
\]
Define $F:\mathcal{P}_W^{ A}\to\mathbb R$ by
\[
F(\pi)
=
\int_{ A}
\frac{d\pi}{d\mu}(a)
\log \frac{d\pi}{d\mu}(a)\,\mu(da),
\]
with the convention $0\log 0=0$. Then $F$ is $W$--strongly convex on
$\mathcal{P}_W^{ A}$ with respect to the $L^2(A,\mu)$ distance between
densities. In particular, for all $\pi,\pi^\prime\in\mathcal{P}_W^{ A}$ such that
$\frac{d\pi}{d\mu}>0$ $\mu$-a.e. and the first variation is well defined,
\[
F(\pi^\prime)
\ge
F(\pi)
+
\left\langle
\frac{\delta F(\pi)}{\delta \pi},
\pi^\prime-\pi
\right\rangle
+
\frac{W}{2}
\left\|
\frac{d\pi^\prime}{d\mu}
-
\frac{d\pi}{d\mu}
\right\|_{L^2(A,\mu)}^2,
\]
where
\[
\frac{\delta F(\pi)}{\delta \pi}
=
1+\log \frac{d\pi}{d\mu}
\]
is the first variation in the sense of Definition
\ref{def:app:flat_derivative}.
\end{proposition}

\begin{proof}
Let
\[
f_\pi:=\frac{d\pi}{d\mu},
\qquad
f_{\pi^\prime}:=\frac{d\pi^\prime}{d\mu}.
\]
By definition of $\mathcal P_W^{ A}$, we have
\[
0\le f_\pi,f_{\pi^\prime}\le W^{-1}
\qquad \mu\text{-a.e.}
\]

Consider $\varphi(x)=x\log x$ on $[0,W^{-1}]$, with the convention
$\varphi(0)=0$. Since
\[
\varphi''(x)=\frac{1}{x}\ge W,
\qquad x\in(0,W^{-1}],
\]
the function $\varphi$ is $W$--strongly convex on $[0,W^{-1}]$. Hence, for
$x,y\in[0,W^{-1}]$ with $x>0$,
\[
\varphi(y)
\ge
\varphi(x)+\varphi'(x)(y-x)+\frac{W}{2}(y-x)^2.
\]
Taking
\[
x=f_\pi(a),
\qquad
y=f_{\pi^\prime}(a),
\]
and integrating over $ A$ gives
\[
F(\pi^\prime)
\ge
F(\pi)
+
\int_{ A}
\left(
1+\log f_\pi(a)
\right)
\left(
f_{\pi^\prime}(a)-f_\pi(a)
\right)
\,\mu(da)
+
\frac{W}{2}
\int_{ A}
\left(
f_{\pi^\prime}(a)-f_\pi(a)
\right)^2
\,\mu(da).
\]
Equivalently,
\[
F(\pi^\prime)
\ge
F(\pi)
+
\left\langle
\frac{\delta F(\pi)}{\delta \pi},
\pi^\prime-\pi
\right\rangle
+
\frac{W}{2}
\left\|
\frac{d\pi^\prime}{d\mu}
-
\frac{d\pi}{d\mu}
\right\|_{L^2(A,\mu)}^2.
\]
Moreover, since both $\pi$ and $\pi^\prime$ are probability measures,
\[
\int_{ A}
\left(
f_{\pi^\prime}(a)-f_\pi(a)
\right)\,\mu(da)
=0,
\]
so the linear term may also be written as
\[
\int_{ A}
\log \frac{d\pi}{d\mu}(a)
\left(
\frac{d\nu}{d\mu}(a)
-
\frac{d\pi}{d\mu}(a)
\right)
\,\mu(da).
\]
This proves the claimed strong convexity inequality.
\end{proof}

Next, let us introduce the KL-logit inequality (for discrete case, refer to Lemma 27~\cite{mei2020global}).
\begin{lemma}\label{lem:app:kl-logit_ineq_affine}
Fix $s\in S$. 
If $\theta,\theta' \in \mathbb{R}^p$ satisfy
\[
\left|
\frac{d\pi_\theta}{d\mu}(\cdot|s)
\right|_{\mathcal B_b( A)}
\le W^{-1},
\qquad
\left|
\frac{d\pi_{\theta'}}{d\mu}(\cdot|s)
\right|_{\mathcal B_b( A)}
\le W^{-1}
\]
for some $W>0$, then
\[
\mathrm{KL}\left(
\pi_\theta(\cdot|s)\,|\,\pi_{\theta'}(\cdot|s)
\right)
\le
\frac{1}{2W}\|\theta-\theta'\|_2^2 .
\]
\end{lemma}

\begin{proof}
For simplicity, write
\[
f_\theta(a|s):=\frac{d\pi_\theta}{d\mu}(a|s),
\qquad
f_{\theta'}(a|s):=\frac{d\pi_{\theta'}}{d\mu}(a|s).
\]
By definition, both densities are strictly positive $\mu$-a.e. Moreover,
by assumption,
\[
0<f_\theta(a|s),f_{\theta'}(a|s)\le W^{-1}
\qquad \mu\text{-a.e.}
\]

Define the negative entropy functional
\[
F(f):=\int_{ A} f(a)\log f(a)\,\mu(da).
\]
By Proposition \ref{prp:app:Strong_convexity_of_negative_entropy}, $F$ is
$W$--strongly convex on the set of densities bounded by $W^{-1}$. Hence,
\[
D_F(f_{\theta'},f_\theta)
\ge
\frac{W}{2}
\|f_{\theta'}-f_\theta\|_{L^2(A,\mu)}^2,
\]
where
\[
D_F(f_{\theta'},f_\theta)
:=
F(f_{\theta'})
-
F(f_\theta)
-
\int_{ A}
\left(1+\log f_\theta(a|s)\right)
\left(f_{\theta'}(a|s)-f_\theta(a|s)\right)
\,\mu(da).
\]
Since both $f_\theta$ and $f_{\theta'}$ are probability densities,
\[
\int_{ A}
\left(f_{\theta'}(a|s)-f_\theta(a|s)\right)\,\mu(da)=0.
\]
Therefore,
\[
D_F(f_{\theta'},f_\theta)
=
\mathrm{KL}\left(
\pi_{\theta'}(\cdot|s)\, |\,\pi_\theta(\cdot|s)
\right).
\]

Using the identity
\[
D_F(f_\theta,f_{\theta'})
+
D_F(f_{\theta'},f_\theta)
=
\int_{ A}
\left(
\log f_\theta(a|s)-\log f_{\theta'}(a|s)
\right)
\left(
f_\theta(a|s)-f_{\theta'}(a|s)
\right)
\,\mu(da),
\]
we obtain
\begin{align*}
&\mathrm{KL}\left(
\pi_\theta(\cdot|s)\,|\,\pi_{\theta'}(\cdot|s)
\right) \\
&\le
\int_{ A}
\left(
\log f_\theta(a|s)-\log f_{\theta'}(a|s)
\right)
\left(
f_\theta(a|s)-f_{\theta'}(a|s)
\right)
\,\mu(da)
-
\frac{W}{2}
\|f_\theta-f_{\theta'}\|_{L^2(A,\mu)}^2 .
\end{align*}
By definition,
\[
\log f_\theta(a|s)-\log f_{\theta'}(a|s)
=
\langle \theta-\theta',g(s,a)\rangle
-
\log Z_\theta(s)
+
\log Z_{\theta'}(s).
\]
The normalizing constants vanish after integration because
\[
\int_{ A}
\left(
f_\theta(a|s)-f_{\theta'}(a|s)
\right)\,\mu(da)=0.
\]
Thus,
\begin{align*}
&\mathrm{KL}\left(
\pi_\theta(\cdot|s)\,|\,\pi_{\theta'}(\cdot|s)
\right) \\
&\le
\int_{ A}
\langle \theta-\theta',g(s,a)\rangle
\left(
f_\theta(a|s)-f_{\theta'}(a|s)
\right)
\,\mu(da)
-
\frac{W}{2}
\|f_\theta-f_{\theta'}\|_{L^2(A,\mu)}^2 .
\end{align*}
By Cauchy's inequality,
\begin{align*}
&\mathrm{KL}\left(
\pi_\theta(\cdot|s)\,|\,\pi_{\theta'}(\cdot|s)
\right) \\
&\le
\|\langle \theta-\theta',g(s,\cdot)\rangle\|_{L^2(A,\mu)}
\|f_\theta-f_{\theta'}\|_{L^2(A,\mu)}
-
\frac{W}{2}
\|f_\theta-f_{\theta'}\|_{L^2(A,\mu)}^2 .
\end{align*}
Using the elementary inequality
\[
xy-\frac{W}{2}y^2\le \frac{x^2}{2W},
\qquad x,y\ge 0,
\]
with
\[
x=\|\langle \theta-\theta',g(s,\cdot)\rangle\|_{L^2(A,\mu)},
\qquad
y=\|f_\theta-f_{\theta'}\|_{L^2(A,\mu)},
\]
we get
\[
\mathrm{KL}\left(
\pi_\theta(\cdot|s)\,|\,\pi_{\theta'}(\cdot|s)
\right)
\le
\frac{1}{2W}
\|\langle \theta-\theta',g(s,\cdot)\rangle\|_{L^2(A,\mu)}^2 .
\]
Finally, by the bound on $g$,
\[
\|\langle \theta-\theta',g(s,\cdot)\rangle\|_{L^2(A,\mu)}
\le
\|\theta-\theta'\|_2.
\]
Therefore,
\[
\mathrm{KL}\left(
\pi_\theta(\cdot|s)\,|\,\pi_{\theta'}(\cdot|s)
\right)
\le
\frac{1}{2W}\|\theta-\theta'\|_2^2 .
\]
\end{proof}

\begin{lemma}\label{lem:app:kl-logit_ineq}
Let Assumptions \ref{asp:app:mdp:max_set_simplex} hold.
Fix $s\in S$. 
If $\theta,\theta'\in\mathbb R^p$ satisfy
\[
\left|
\frac{d\pi_\theta}{d\mu}(\cdot|s)
\right|_{\mathcal B_b( A)}
\le W^{-1},
\qquad
\left|
\frac{d\pi_{\theta'}}{d\mu}(\cdot|s)
\right|_{\mathcal B_b( A)}
\le W^{-1}
\]
for some $W>0$, then for every $c\in\mathbb R$,
\[
\mathrm{KL}\left(
\pi_\theta(\cdot|s)\,|\,\pi_{\theta'}(\cdot|s)
\right)
\le
\frac{1}{2W}
\left\|
\theta-\theta'-c\mathbf 1
\right\|_2^2 .
\]
\end{lemma}

\begin{proof}
Write
\[
f_\theta(a|s):=\frac{d\pi_\theta}{d\mu}(a|s),
\qquad
f_{\theta'}(a|s):=\frac{d\pi_{\theta'}}{d\mu}(a|s).
\]
By the logit parametrization, both densities are strictly positive
$\mu$-a.e. Moreover, by assumption,
\[
0<f_\theta(a|s),f_{\theta'}(a|s)\le W^{-1},
\qquad \mu\text{-a.e.}
\]

Let
\[
F(f):=\int_{ A} f(a)\log f(a)\,\mu(da)
\]
be the negative entropy functional. By Proposition
\ref{prp:app:Strong_convexity_of_negative_entropy}, $F$ is
$W$--strongly convex on densities bounded by $W^{-1}$. Hence,
\[
F(f_{\theta'})
\ge
F(f_\theta)
+
\int_{ A}
\log f_\theta(a|s)
\left(
f_{\theta'}(a|s)-f_\theta(a|s)
\right)\,\mu(da)
+
\frac{W}{2}
\left\|
f_{\theta'}-f_\theta
\right\|_{L^2(A,\mu)}^2,
\]
where we used
\[
\int_{ A}
\left(
f_{\theta'}(a|s)-f_\theta(a|s)
\right)\,\mu(da)=0
\]
to remove the constant term in the first variation.

Therefore,
\begin{align*}
&\mathrm{KL}\left(
\pi_\theta(\cdot|s)\,|\,\pi_{\theta'}(\cdot|s)
\right) \\
&=
\int_{ A}
\log f_\theta(a|s) f_\theta(a|s)\,\mu(da)
-
\int_{ A}
\log f_{\theta'}(a|s) f_\theta(a|s)\,\mu(da) \\
&\le
\int_{ A}
\left(
\log f_\theta(a|s)-\log f_{\theta'}(a|s)
\right)
\left(
f_\theta(a|s)-f_{\theta'}(a|s)
\right)\,\mu(da) \\
&\qquad
-
\frac{W}{2}
\left\|
f_{\theta'}-f_\theta
\right\|_{L^2(A,\mu)}^2 .
\end{align*}

By the logit parametrization,
\[
\log f_\theta(a|s)-\log f_{\theta'}(a|s)
=
\langle \theta-\theta',g(s,a)\rangle
-\log Z_\theta(s)+\log Z_{\theta'}(s).
\]
Assumption \ref{asp:app:mdp:max_set_simplex} gives
\[
\langle \mathbf 1,g(s,a)\rangle=1.
\]
Thus, for any $c\in\mathbb R$,
\[
\langle \theta-\theta',g(s,a)\rangle
=
\langle \theta-\theta'-c\mathbf 1,g(s,a)\rangle
+
c.
\]
Since both $f_\theta(\cdot|s)$ and $f_{\theta'}(\cdot|s)$ are probability
densities,
\[
\int_{ A}
\left(
f_\theta(a|s)-f_{\theta'}(a|s)
\right)\,\mu(da)=0.
\]
Hence the constant term
\[
c-\log Z_\theta(s)+\log Z_{\theta'}(s)
\]
vanishes after integration against
$f_\theta(\cdot|s)-f_{\theta'}(\cdot|s)$. Therefore,
\begin{align*}
&\mathrm{KL}\left(
\pi_\theta(\cdot|s)\,|\,\pi_{\theta'}(\cdot|s)
\right) \\
&\le
\int_{ A}
\langle \theta-\theta'-c\mathbf 1,g(s,a)\rangle
\left(
f_\theta(a|s)-f_{\theta'}(a|s)
\right)\,\mu(da) \\
&\qquad
-
\frac{W}{2}
\left\|
f_{\theta'}-f_\theta
\right\|_{L^2(A,\mu)}^2 .
\end{align*}
By Cauchy's inequality,
\begin{align*}
&\mathrm{KL}\left(
\pi_\theta(\cdot|s)\,|\,\pi_{\theta'}(\cdot|s)
\right) \\
&\le
\left\|
\langle \theta-\theta'-c\mathbf 1,g(s,\cdot)\rangle
\right\|_{L^2(A,\mu)}
\left\|
f_\theta-f_{\theta'}
\right\|_{L^2(A,\mu)}
-
\frac{W}{2}
\left\|
f_\theta-f_{\theta'}
\right\|_{L^2(A,\mu)}^2 .
\end{align*}
Using
\[
xy-\frac{W}{2}y^2\le \frac{x^2}{2W},
\qquad x,y\ge 0,
\]
we get
\[
\mathrm{KL}\left(
\pi_\theta(\cdot|s)\,|\,\pi_{\theta'}(\cdot|s)
\right)
\le
\frac{1}{2W}
\left\|
\langle \theta-\theta'-c\mathbf 1,g(s,\cdot)\rangle
\right\|_{L^2(A,\mu)}^2 .
\]
Finally, by the bound on $g$,
\[
\left\|
\langle \theta-\theta'-c\mathbf 1,g(s,\cdot)\rangle
\right\|_{L^2(A,\mu)}
\le
\left\|
\theta-\theta'-c\mathbf 1
\right\|_2 .
\]
Therefore,
\[
\mathrm{KL}\left(
\pi_\theta(\cdot|s)\,|\,\pi_{\theta'}(\cdot|s)
\right)
\le
\frac{1}{2W}
\left\|
\theta-\theta'-c\mathbf 1
\right\|_2^2 .
\]
\end{proof}

\begin{lemma}[Concavity of variance]\label{lem:app:concavity_of_var}
Let $P=qP_1+(1-q)P_2$ with $q\in[0,1]$, where $P_1$ and $P_2$ are
probability measures. 
If $X\in L^2(P_1)\cap L^2(P_2)$, then
\[
\operatorname{Var}_P(X)
\ge
q\operatorname{Var}_{P_1}(X)
+
(1-q)\operatorname{Var}_{P_2}(X).
\]
\end{lemma}

\begin{proof}
Let $\bar{\mu}_i=\mathbb E_{P_i}[X]$ for $i=1,2$. Since
$P=qP_1+(1-q)P_2$,
\[
\mathbb E_P[X]=q\bar{\mu}_1+(1-q)\bar{\mu}_2,
\qquad
\mathbb E_P[X^2]
=
q\mathbb E_{P_1}[X^2]+(1-q)\mathbb E_{P_2}[X^2].
\]
Therefore,
\[
\begin{aligned}
\operatorname{Var}_P(X)
&=
q\operatorname{Var}_{P_1}(X)
+
(1-q)\operatorname{Var}_{P_2}(X)
+
q(1-q)(\bar{\mu}_1-\bar{\mu}_2)^2 \\
&\ge
q\operatorname{Var}_{P_1}(X)
+
(1-q)\operatorname{Var}_{P_2}(X).
\end{aligned}
\]
\end{proof}

For any $\pi \in \Pi_\mu$, define
\[
G^{\pi}(s) := \int_A \left(g(s,a) - \int_A g(s, a^\prime) \pi ( da^\prime | s) \right)\left(g(s,a) - \int_A g(s, a^\prime) \pi (da^\prime | s) \right)^{\top}\pi(da|s).
\]

\begin{lemma}\label{lemapp:mdp:lower_bound_log_density}
Let $\theta^\prime \in \mathbb{R}^p$.
If there exists a constant $R > 0$ such that $\left|\log \frac{d\pi}{d\mu}(a|s)\right|_{B_b({S \times A})} \leq  R$, then there exists 
\[
\lambda =  e^{-R - 2\|\theta^\prime\|_2} \int_{S}  \lambda_{\min} \left(G^{\pi_{\theta^\prime}}(s) \right)d_\rho^{\pi^*}(ds)
\]
such that $\int_{S} \lambda_{\min} (G^\pi(s))d_\rho^{\pi^*}(ds) \geq \lambda$.
\end{lemma}
\begin{proof}
For all $s \in S, a \in A$,
\begin{align*}
&\left|\log\frac{d \pi_{\theta^\prime}}{d \mu}(a|s)\right|\\
& = \left| \langle \theta^\prime, g(s,a)\rangle - \log Z_{\pi_{\theta^\prime}}(s)\right| \\
& \leq \sup_{s,a}\|g(s,a)\|_2\|\theta^\prime\|_2 + \left|\log \int_{A} e^{\langle \theta^\prime,g(s,a)\rangle}\mu(da)\right|\\
& \leq 2\|\theta^\prime\|_2,
\end{align*}
so $\left|\frac{d \pi_{\theta^\prime}}{d \mu}\right|_{B_b(S \times A)} \leq e^{2\|\theta^\prime\|_2}$.
Since $\min_{s \in S, a\in A} \frac{d\pi}{d\pi_{\theta^\prime}}(a|s) = \min_{s \in S, a\in A} \frac{d\pi}{d\mu} \frac{d\mu}{d \pi_{\theta^\prime}}(a|s) \geq e^{-R - 2\|\theta^\prime\|_2}  > 0$, then we know $e^{-R - 2\|\theta^\prime\|_2}  \leq \int_{A} \frac{d\pi}{d \pi_{\theta^\prime}}(a|s)\pi_{\theta^\prime}(da) = 1$, and
\[
\pi(da|s) = e^{-R - 2\|\theta^\prime\|_2}  \pi_{\theta_{0}}(da) + \left(1-e^{-R - 2\|\theta^\prime\|_2} \right) \bar{\pi}(da|s),
\]
where $\bar{\pi}(da|s) = \frac{\pi(da|s)- e^{-R - 2\|\theta^\prime\|_2} \pi_{\theta^\prime}(da)}{1-e^{-R - 2\|\theta^\prime\|_2} }$ determines a valid probability measure for fixed $s$. 

Hence by Lemma \ref{lem:app:concavity_of_var},
\begin{align*} 
&\mathrm{Var}_{a\sim\pi(\cdot | s)}(v^{\top}g(s,a))\\
& \geq e^{-R - 2\|\theta^\prime\|_2} \mathrm{Var}_{a\sim\pi_{\theta^\prime}}(v^{\top}g(s,a))  + \left(1-e^{-R - 2\|\theta^\prime\|_2} \right)\mathrm{Var}_{a\sim\bar{\pi}(\cdot | s)}(v^{\top}g(s,a)) \\
& \geq e^{-R - 2\|\theta^\prime\|_2} \mathrm{Var}_{a\sim\pi_{\theta^\prime}}(v^{\top}g(s,a))\\
& =  q e^{-2\|\theta^\prime\|_2} v^{\top} G^{\pi_{\theta^\prime}}(s) v\\
& \geq e^{-R - 2\|\theta^\prime\|_2}  \lambda_{\min} \left(G^{\pi_{\theta^\prime}}(s)\right) ,
\end{align*}
where in the first equality we use the Definition \ref{def:app:mdp:FIM} of FIM,
{\small
\begin{align*}
&\mathrm{Var}_{a\sim\pi_{\theta^\prime}}(v^{\top}g(s,a)) \\
&=\int_A \left(v^{\top}g(s,a) - \int_A v^{\top}g(s, a^\prime) \pi_{\theta^\prime} ( da^\prime | s) \right)\left(v^{\top}g(s,a) - \int_A v^{\top}g(s, a^\prime) \pi_{\theta^\prime} (da^\prime | s) \right)\pi_{\theta^\prime}(da|s)\\
& = v^{\top} G^{\pi_{\theta^\prime}}(s) v,
\end{align*}
}
which implies that,
\begin{align*}
\lambda_{\min}(G^\pi)(s)& = \min_{\|v\|_2 = 1} v^\top G^\pi(s) v \quad \\
& = \min_{\|v\|_2 = 1} \mathrm{Var}_{a\sim\pi(\cdot | s)}(v^{\top}g(s,a))\\
& \geq e^{-R - 2\|\theta^\prime\|_2} \lambda_{\min} \left(G^{\pi_{\theta^\prime}}(s)\right)\\
& = e^{-R - 2\|\theta^\prime\|_2}  \lambda_{\min} \left(G^{\pi_{\theta^\prime}}(s)\right),
\end{align*}
then there exists 
\[
\lambda = e^{-R - 2\|\theta^\prime\|_2} \int_{A} \lambda_{\min} \left(G^{\pi_{\theta^\prime}}(s) \right)d_\rho^{\pi_{\theta^*}}(ds)
\]
such that $\int_{A}\lambda_{\min} (G^\pi(s))d_\rho^{\pi^*}(ds) \geq \lambda$.
\end{proof}

\begin{lemma}\label{lem:app:mdp:lower_bound_log_density_simplex}
Let $\theta^\prime \in \mathbb{R}^p$.
If there exists a constant $R > 0$ such that $\left|\log \frac{d\pi}{d\mu}(a|s)\right|_{B_b({S \times A})} \leq  R$, then there exists 
\[
\lambda =  e^{-R - 2\|\theta^\prime\|_2} \int_{S}  \lambda_{\min}\left(\int g(s,a)g^\top(s,a)\pi_{\theta^\prime} (da)\right)d^{\pi^*}_\rho(ds)
\]
such that $\int\lambda_{\min} \left(\int g(s,a)g^\top(s,a)\pi (da | s)\right))d_\rho^{\pi_{\theta^*}}(ds) \geq \lambda$.
\end{lemma}
\begin{proof}

Also for all $s \in S, a \in A$,
\begin{align*}
&\left|\log\frac{d \pi_{\theta^\prime}}{d \mu}(a|s)\right|\\
& = \left| \langle \theta^\prime, g(s,a)\rangle - \log Z_{\pi_{\theta^\prime}}(s)\right| \\
& \leq \sup_{s,a}\|g(s,a)\|_2\|\theta^\prime\|_2 + \left|\log \int e^{\langle \theta^\prime,g(s,a)\rangle}\mu(da)\right|\\
& \leq 2\|\theta^\prime\|_2,
\end{align*}
so $\left|\frac{d \pi_{\theta^\prime}}{d \mu}\right|_{B_b(S \times A)} \leq e^{2\|\theta^\prime\|_2}$.

Since $\min_{s \in S, a\in A} \frac{d\pi}{d\pi_{\theta^\prime}}(a|s) = \min_{s \in S, a\in A} \frac{d\pi}{d\mu} \frac{d\mu}{d \pi_{\theta^\prime}}(a|s) \geq e^{-R - 2\|\theta^\prime\|_2} >0$, then we know $e^{-R - 2\|\theta^\prime\|_2} \leq \int_{A} \frac{d\pi}{d \pi_{\theta^\prime}}(a|s)\pi_{\theta^\prime}(da) = 1$, and
\[
\pi(da|s) = e^{-R - 2\|\theta^\prime\|_2} \pi_{\theta_{0}}(da) + \left(1-e^{-R - 2\|\theta^\prime\|_2}\right) \bar{\pi}(da|s),
\]
where $\bar{\pi}(da|s) = \frac{\pi(da|s)- e^{-R - 2\|\theta^\prime\|_2}\pi_{\theta^\prime}(da)}{1-e^{-R - 2\|\theta^\prime\|_2}}$ determines a valid probability measure for fixed $s$. 

Hence by Lemma \ref{lem:app:concavity_of_var},
\begin{align*}
&\lambda_{\min}\left(\int_{A} g(s,a)g^\top(s,a)\pi(da|s)\right)\\
& = \min_{\|v\|_2 = 1} v^\top \left(\int_{A} g(s,a)g^\top(s,a)\pi(da|s)\right) v\\
& = \min_{\|v\|_2 = 1} v^\top \left(\int_{A} g(s,a)g^\top(s,a) \left(e^{-R - 2\|\theta^\prime\|_2}\pi_{\theta^\prime}(da | s) + \left(1-e^{-R - 2\|\theta^\prime\|_2}\right) \bar{\pi}(da|s)\right)\right) v\\
& \geq e^{-R - 2\|\theta^\prime\|_2} \min_{\|v\| = 1} v^\top \left(\int_{A} g(s,a)g^\top(s,a) \pi_{\theta^\prime}(da | s)\right) v \\
& = e^{-R - 2\|\theta^\prime\|_2} \lambda_{\min}\left(\int_{A} g(s,a)g^\top(s,a)\pi_{\theta^\prime} (da | s)\right)>0.
\end{align*}
\end{proof}


\begin{proposition}\label{prp:app:mdp:singularity_of_fim_mdp_simplex}
Let Assumption \ref{asp:app:mdp:max_set_affine} hold.
For all $\theta \in \mathbb{R}^p$, $\lambda_{\min} (G^{\pi_\theta}(s)) = 0$ with eigenvector $\textbf{1} \in \mathbb{R}^p$.
\end{proposition}
\begin{proof}
Define $\bar{g}_\theta (s):= \int_A g(s,a) \pi_\theta(da|s)$, and note that $\bar{g}_\theta^\top (s)\textbf{1} = \int g(s,a)^\top \textbf{1}\pi_\theta(da|s) =1$, then we have
\begin{align*}
G^{\pi_\theta} (s)\textbf{1} &= \int_A \nabla\log\frac{d \pi_{\theta}}{d \mu}(a|s)\left(\nabla\log\frac{d \pi_{\theta}}{d \mu}(a|s)\right)^{\top} \textbf{1}\pi_\theta(da|s)\\
& = \int_A \left(g(s,a) - \bar{g}_\theta(s)\right)\left(g(s,a) - \bar{g}_\theta(s)\right)^{\top} \textbf{1}\pi_\theta(da|s)\\
& = \mathbf{0}.
\end{align*}
\end{proof}

Next, we introduce the following interlacing theorem of eigenvalues of a real symmetric matrix perturbed by a rank 1 matrix (see Section 5 in~\cite{4d2f28f8-3cb9-3c33-89f8-2a7e91475e75}).
\begin{theorem}[Eigenvalue Interlacing Theorem]\label{thm:app:bandit:Eigenvalue_Interlacing}
Let B be a real symmetric matrix. Define:
\[
A = B + xx^T
\]
If the n eigenvalues of A are $\eta_1 \ge \dots \ge \eta_n$, and the n eigenvalues of B are $\lambda_1 \ge \dots \ge \lambda_n$, then these are interlaced as: $\eta_1 \ge \lambda_1 \ge \dots \ge \eta_i \ge \lambda_i \ge \eta_{i+1} \dots \eta_n \ge \lambda_n$.
\end{theorem}

The following Corollary \ref{cor:app:mdp:spectrum_of_fim} is similar to Lemma 23 in~\cite{mei2020global}.

\begin{corollary}\label{cor:app:mdp:spectrum_of_fim}
Let Assumption \ref{asp:app:mdp:max_set_affine} hold. Then for fixed $s \in \mathcal{S}$, for any vector $y, \theta \in \mathbb{R}^p$, then 
\[
{\left\|G^{\pi_\theta}(s)\left(y-\frac{y^{\top} \mathbf{1}}{p} \cdot \mathbf{1}\right)\right\|_2 \geq \lambda_{\min}\left(\int_{A} g(s,a)g^\top(s, a)\pi_\theta(da|s)\right)\cdot\left\|y-\frac{y^{\top} \mathbf{1}}{p} \cdot \mathbf{1}\right\|_2} .
\]
\end{corollary}
\begin{proof}
Let $\eta_i, i=2,3, \ldots, p$ be the eigenvalue of matrix $\int_{A} g(s,a)g^\top(s, a)\pi_\theta(da|s)$ in the following order
\[
\eta_1 \geq \eta_2 \geq \cdots \geq \eta_p,
\]
and denote the eigenvalues of $G^{\pi_\theta}(s)$ as
\[
\lambda_1 \geq \lambda_2 \geq \cdots \geq \lambda_p .
\]
Then Theorem \ref{thm:app:bandit:Eigenvalue_Interlacing} indicates that,
\[
\eta_{i+1} \leq \lambda_{i} \leq  \eta_i, i=1,2,3 \ldots, p - 1.
\]
Since $G^{\pi_\theta}(s)$ is real symmetric, its eigenvectors $\left\{\frac{\mathbf 1}{\sqrt{p}}, v_{p-1}, \ldots, v_1\right\}$ are orthonormal. For any vector $y$, $y$ can be written as linear combination of eigenvectors of $G^{\pi_\theta}(s)$,
\[
\begin{aligned}
y & =a_1 \cdot \frac{\mathbf 1}{\sqrt{p}}+a_{2}v_{p-1}+\cdots+a_p v_1 \\
& =\frac{y^{\top} \mathbf{1}}{p} \cdot \mathbf{1}+a_2 v_{p-1}+\cdots+a_p v_1.
\end{aligned}
\]
The last equation is because the representation is unique, and
\[
a_p=y^{\top} \frac{\mathbf{1}}{\sqrt{p}}=\frac{y^{\top} \mathbf{1}}{\sqrt{p}}.
\]
Denote
\[
y^{\prime}=y-\frac{y^{\top} \mathbf{1}}{K} \cdot \mathbf{1}=a_{2} v_{p-1}+\cdots+a_p v_1.
\]

We have
\[
\left\|y^{\prime}\right\|_2^2=a_2^2+\cdots+a_{p}^2.
\]
Therefore because of $\lambda_{p-1}\geq \eta_{p}= \lambda_{\min}\left(\int_{A} g(s,a)g^\top(s,a)\pi_\theta(da|s)\right)$,
\[
\begin{aligned}
\left\|G^{\pi_\theta}(s) y^{\prime}\right\|_2 & =\left(a_{2}^2 \lambda_{p-1}^2+\cdots+a_p^2 \lambda_1^2\right)^{\frac{1}{2}} \\
& \geq\left(\left(a_2^2+\cdots+a_{p}^2\right) \cdot \lambda_{p-1}^2\right)^{\frac{1}{2}} \\
& =\lambda_{p-1}\cdot\left\|y^{\prime}\right\|_2 \\
& \geq \lambda_{\min}\left(\int_{A} g(s,a)g^\top(s,a)\pi_\theta(da|s)\right) \cdot\left\|y^{\prime}\right\|_2 .
\end{aligned}
\]
\end{proof}

\subsubsection{Proof of Proposition \ref{prp:app:mdp:Gradient_of_the_Objective}}
\label{sec:add_proof}

\begin{proof}
Let $\theta, \theta' \in \mathbb{R}^p$ and define $\theta^{\varepsilon}=\theta+\varepsilon (\theta'-\theta)$ for $\varepsilon\in[0,1]$. 
\begin{align*}
\nabla V_\tau^{\pi_\theta}(\rho) &= \lim_{\varepsilon \rightarrow 0}\frac{V_\tau^{\pi_{\theta^\epsilon}}(\rho) - V_\tau^{\pi_\theta}(\rho)}{\varepsilon}
\end{align*}
Noting that  
\begin{equation}\label{eq:Vtau_diff_quot}
\frac{V_\tau^{\pi_{\theta^\epsilon}}(\rho) - V_\tau^{\pi_\theta}(\rho)}{\varepsilon} = \int_{S} \frac{V^{\pi_{\theta^{\varepsilon}}}_{\tau}(s)-V^{\pi_{\theta}}_{\tau}(s)}{\varepsilon}\rho(ds)\,,
\end{equation}
we first study the difference quotient
$
H^{\varepsilon}(s)=\frac{V^{\pi_{\theta^{\varepsilon}}}_{\tau}(s)-V^{\pi_{\theta}}_{\tau}(s)}{\varepsilon}\,.
$
By \eqref{eq:intro:policy_Bellman}, we have
\begin{equation}\label{eq:first_iterate_value}
\begin{split}
H^{\varepsilon}(s)&=\frac{1}{\varepsilon}\left[\int_{A}\left(c(s,a)+\gamma\int_{S}P(ds'|s,a)V^{\pi_{\theta^{\varepsilon}}}_{\tau}(s')+\tau\log\frac{d\pi_{\theta^{\varepsilon}}}{d\mu}(a|s)\right)\pi_{\theta^{\varepsilon}}(da|s)\right.\\
&\left.\qquad-\int_{A}\left(c(s,a)+\gamma\int_{S}P(ds'|s,a)V^{\pi_{\theta}}_{\tau}(s')+\tau\log\frac{d\pi_{\theta}}{d\mu}(a|s)\right)\pi_{\theta}(da|s)\right]\\
&=I^{\varepsilon}_1(s) + I^{\varepsilon}_2(s)+I^{\varepsilon}_3(s) +I^{\varepsilon}_4(s)+I^{\varepsilon}_5(s)\,,
\end{split}
\end{equation}
where
\begin{align*}
I^{\varepsilon}_1(s)&:=\int_{A}c(s,a)\frac{\pi_{\theta^{\varepsilon}}(da|s)-\pi_{\theta}(da|s)}{\varepsilon}\,,\\
I^{\varepsilon}_2(s)&:=\gamma\int_{A}\int_{S}\frac{V^{\pi_{\theta^{\varepsilon}}}_{\tau}(s')-V^{\pi_{\theta}}_{\tau}(s')}{\varepsilon}P(ds'|s,a)\pi_{\theta}(da|s)=\gamma\int_{S} H^{\varepsilon}(s')P_{\pi_{\theta}}(ds'|s)\,,\\
I^{\varepsilon}_3(s)&:=\gamma\int_{A}\int_{S}V^{\pi_{\theta^{\varepsilon}}}_{\tau}(s')P(ds'|s,a)\frac{\pi_{\theta^{\varepsilon}}(da|s)-\pi_{\theta}(da|s)}{\varepsilon}\,,\\
I^{\varepsilon}_4(s)&:=+\tau \int_{A}\frac{\log\frac{d\pi_{\theta^{\varepsilon}}}{d\mu}(a|s) -\log\frac{d\pi_{\theta}}{d\mu}(a|s) }{\varepsilon}\pi_{\theta^{\varepsilon}}(da|s)\,\\
I^{\varepsilon}_5(s)&:=+\tau\int_{A}\log\frac{d\pi_{\theta}}{d\mu}(a|s)\frac{\pi_{\theta^{\varepsilon}}(da|s)-\pi_{\theta}(da|s)}{\varepsilon}\,. 
\end{align*}
Iterating \eqref{eq:first_iterate_value}, we obtain
\begin{align*}
H^{\varepsilon}(s) &= I^{\varepsilon}_1(s) + I^{\varepsilon}_3(s)+I^{\varepsilon}_4(s)+I^{\varepsilon}_5(s) +\gamma  \int_{S}H^{\varepsilon}(s')P_{\pi_{\theta}}(ds'|s)\\
&=\frac{1}{1-\gamma} \int_{S} (I^{\varepsilon}_1(s')+I^{\varepsilon}_3(s')+I^{\varepsilon}_4(s')+I^{\varepsilon}_5(s'))d^{\pi_{\theta}}(ds'|s)\,,
\end{align*}
and hence
\[
\frac{V_\tau^{\pi_{\theta^\epsilon}}(\rho) - V_\tau^{\pi_\theta}(\rho)}{\varepsilon} =\frac{1}{1-\gamma} \int_{S} (I^{\varepsilon}_1(s)+I^{\varepsilon}_3(s)+I^{\varepsilon}_4(s)+I^{\varepsilon}_5(s))d^{\pi_{\theta}}_{\rho}(ds)\,.
\]
\textbf{Step 1}: To prove \begin{equation}\label{eq:limit_policy}
\underset{\varepsilon \in [0,1]}{\lim_{\varepsilon\rightarrow 0}}|\pi_{\theta^{\varepsilon}}-\pi_{\theta}|_{b\mathcal{K}(A|S)}=\underset{\varepsilon \in [0,1]}{\lim_{\varepsilon\rightarrow 0}}\sup_{s\in S}\int_{A}\left|\frac{d\pi_{\theta^{\varepsilon}}}{d\mu}(a|s)-\frac{d\pi_{\theta}}{d\mu}(a|s)\right|\mu(da)=0\,.
\end{equation}

For convenience, we introduce the unnormalized policy $\tilde{\pi}: \mathcal{P}(\mathbb{R}^d) \rightarrow b\mathcal{K}_{\mu}(A|S)$ given by
\[
\tilde{\pi}(\theta)(da|s)=\tilde{\pi}_{\theta}(da|s)=\exp\left(\langle \theta, g(s,a)\rangle\right)\mu(da)\,.
\]
For all $(s,a)\in S\times A$, we have
\begin{align*}
& \frac{d\pi_{\theta^{\varepsilon}}}{d\mu}(a|s)-\frac{d\pi_{\theta}}{d\mu}(a|s) \\
& =\frac{\frac{d\tilde{\pi}_{\theta^{\varepsilon}}}{d\mu}(a|s)}{\tilde{\pi}_{\theta^{\varepsilon}}(A|s)}-\frac{\frac{d\tilde{\pi}_{\theta}}{d\mu}(a|s)}{\tilde{\pi}_{\theta}(A|s)}=\frac{\frac{d\tilde{\pi}_{\theta^{\varepsilon}}}{d\mu}(a|s)\tilde{\pi}_{\theta}(A|s)}{\tilde{\pi}_{\theta^{\varepsilon}}(A|s)\tilde{\pi}_{\theta}(A|s)}-\frac{\frac{d\tilde{\pi}_{\theta}}{d\mu}(a|s)\tilde{\pi}_{\theta^{\varepsilon}}(A|s)}{\tilde{\pi}_{\theta^{\varepsilon}}(A|s)\tilde{\pi}_{\theta}(A|s)}\\
& =\frac{\frac{d\tilde{\pi}_{\theta^{\varepsilon}}}{d\mu}(a|s)-\frac{d\tilde{\pi}_{\theta}}{d\mu}(a|s)}{\tilde{\pi}_{\theta}(A|s)}\frac{\tilde{\pi}_{\theta}(A|s)}{\tilde{\pi}_{\theta^{\varepsilon}}(A|s)}+\frac{\frac{d\tilde{\pi}_{\theta}}{d\mu}(a|s)}{\tilde{\pi}_{\theta^{\varepsilon}}(A|s)}\frac{\tilde{\pi}_{\theta}(A|s)-\tilde{\pi}_{\theta^{\varepsilon}}(A|s)}{\tilde{\pi}_{\theta}(A|s)}\\
& =\left[\frac{\frac{d\tilde{\pi}_{\theta^{\varepsilon}}}{d\mu}(a|s)-\frac{d\tilde{\pi}_{\theta}}{d\mu}(a|s)}{\tilde{\pi}_{\theta}(A|s)}+\frac{d\pi_{\theta}}{d\mu}(a|s)\frac{\tilde{\pi}_{\theta}(A|s)-\tilde{\pi}_{\theta^{\varepsilon}}(A|s)}{\tilde{\pi}_{\theta}(A|s)}\right]\frac{\tilde{\pi}_{\theta}(A|s)}{\tilde{\pi}_{\theta^{\varepsilon}}(A|s)}\\
& =\left[\frac{\frac{d\tilde{\pi}_{\theta^{\varepsilon}}}{d\mu}(a|s)-\frac{d\tilde{\pi}_{\theta}}{d\mu}(a|s)}{\tilde{\pi}_{\theta}(A|s)}+\frac{d\pi_{\theta}}{d\mu}(a|s)\int_{A}\left(\frac{\frac{d\tilde{\pi}_{\theta^{\varepsilon}}}{d\mu}(a|s)-\frac{d\tilde{\pi}_{\theta}}{d\mu}(a|s)}{\tilde{\pi}_{\theta}(A|s)}\right)\mu(da)\right]\frac{\tilde{\pi}_{\theta}(A|s)}{\tilde{\pi}_{\theta^{\varepsilon}}(A|s)}\,.
\end{align*}
Simple manipulation yields
\[
\frac{d\tilde{\pi}_{\theta^{\varepsilon}}}{d\mu}(a|s)-\frac{d\tilde{\pi}_{\theta}}{d\mu}(a|s) =\exp\left(\langle\theta, g(s,a)\rangle\right)\left( \exp\left(\varepsilon\langle\theta^\prime - \theta, g(s,a)\rangle\right)-1\right)\,.
\]
Taylor expanding the exponential function, we find
\begin{align*}
& \exp\left(\varepsilon\langle\theta^\prime - \theta, g(a)\rangle\right)-1\\
& = \varepsilon\langle\theta^\prime - \theta, g(s,a)\rangle + \varepsilon^2\sum_{n=2}^\infty\varepsilon^{n-2} \frac{\varepsilon\langle\theta^\prime - \theta, g(s,a)\rangle^n}{n!}\,,
\end{align*}
and hence
\begin{equation}
\begin{aligned}\label{ineq:ratio_diff_pol_eps}
\varepsilon^{-1}\frac{\frac{d\tilde{\pi}_{\theta^{\varepsilon}}}{d\mu}(a|s)-\frac{d\tilde{\pi}_{\theta}}{d\mu}(a|s)}{\tilde{\pi}_{\theta}(A|s)}&= \frac{d\pi_{\theta}}{d\mu}(a|s)\langle\theta^\prime - \theta, g(s,a)\rangle\\
&\quad +\varepsilon \frac{d\pi_{\theta}}{d\mu}(a|s) \sum_{n=2}^\infty \varepsilon^{n-2}\frac{\varepsilon\langle\theta^\prime - \theta, g(s,a)\rangle^n}{n!}\,.
\end{aligned}
\end{equation}
Thus, using $\pi_{\theta}(A|s)=1$, we obtain
\begin{gather*}
\int_{A}\left|\varepsilon^{-1}\frac{\frac{d\tilde{\pi}_{\theta^{\varepsilon}}}{d\mu}(a|s)-\frac{d\tilde{\pi}_{\theta}}{d\mu}(a|s)}{\tilde{\pi}_{\theta}(A|s)}- \frac{d\pi_{\theta}}{d\mu}(a|s)\langle\theta^\prime - \theta, g(a)\rangle\right|\mu(da)\notag\\
\le  \varepsilon\exp\left(\|\theta'-\theta\|_2\right)
\end{gather*}
and 
\[
\int_{A}\varepsilon^{-1}\left|\frac{\frac{d\tilde{\pi}_{\theta^{\varepsilon}}}{d\mu}(a|s)-\frac{d\tilde{\pi}_{\theta}}{d\mu}(a|s)}{\tilde{\pi}_{\theta}(A|s)}\right|\mu(da) \le \|\theta'-\theta\|_{2}+\varepsilon \exp\left(\|\theta'-\theta\|_{2}\right).
\]
The dominated convergence theorem implies that 
\begin{equation*}
\underset{\varepsilon \in [0,1]}{\lim_{\varepsilon\rightarrow 0}}|\pi_{\theta^{\varepsilon}}-\pi_{\theta}|_{b\mathcal{K}(A|S)}=\underset{\varepsilon \in [0,1]}{\lim_{\varepsilon\rightarrow 0}}\sup_{s\in S}\int_{A}\left|\frac{d\pi_{\theta^{\varepsilon}}}{d\mu}(a|s)-\frac{d\pi_{\theta}}{d\mu}(a|s)\right|\mu(da)=0\,,
\end{equation*}
\textbf{Step 2}: We now will pass to the limit as $\varepsilon\rightarrow 0$ in Eq.(\ref{eq:Vtau_diff_quot}). 
Let us begin with the $I_4^{\varepsilon}$-term. Recalling \eqref{eq:limit_policy}, we have
\begin{equation}\label{eq:limit_policy_again}
\underset{\varepsilon \in [0,1]}{\lim_{\varepsilon\rightarrow 0}}|\pi_{\theta^{\varepsilon}}-\pi_{\theta}|_{b\mathcal{K}(A|S)}=\underset{\varepsilon \in [0,1]}{\lim_{\varepsilon\rightarrow 0}}\sup_{s\in S}\int_{A}\left|\frac{d\pi_{\theta^{\varepsilon}}}{d\mu}(a|s)-\frac{d\pi_{\theta}}{d\mu}(a|s)\right|\mu(da)=0\,.
\end{equation}
Since $\pi_\theta \in \Pi_\mu$, there exists a constant $\tilde{R}_\theta >0$ such that $\left|\frac{d \pi_{\theta}}{d \mu}\right|_{B_b(S \times A)} \leq \tilde{R}_\theta $
\[
\frac{d\pi_{\theta}}{d\mu}(a|s) \geq e^{-\tilde{R}_\theta}>0,
\]
there is an $\varepsilon_0\in (0,1]$ such that  for all $\varepsilon<\varepsilon_0$, $s\in S$, and $\mu-a.e. \,a\in A$,
\[
\left|\frac{\frac{d\pi_{\theta^{\varepsilon}}}{d\mu}(a|s)-\frac{d\pi_{\theta}}{d\mu}(a|s)}{\frac{d\pi_{\theta}}{d\mu}(a|s)}\right|<\frac{1}{2}\,.
\]
Taylor expanding the logarithm, we get 
\begin{gather*}
\log\frac{d\pi_{\theta^{\varepsilon}}}{d\mu}(a|s) -\log\frac{d\pi_{\theta}}{d\mu}(a|s) =\log\left(1+\frac{\frac{d\pi_{\theta^{\varepsilon}}}{d\mu}(a|s)-\frac{d\pi_{\theta}}{d\mu}(a|s)}{\frac{d\pi_{\theta}}{d\mu}(a|s)}\right)\notag\\
=\frac{1}{\frac{d\pi_{\theta}}{d\mu}(a|s)}\left(\frac{d\pi_{\theta^{\varepsilon}}}{d\mu}(a|s)-\frac{d\pi_{\theta}}{d\mu}(a|s)\right)+\sum_{n=2}^{\infty}(-1)^{n+1}\frac{\left(\frac{1}{\frac{d\pi_{\theta}}{d\mu}(a|s)}\left(\frac{d\pi_{\theta^{\varepsilon}}}{d\mu}(a|s)-\frac{d\pi_{\theta}}{d\mu}(a|s)\right)\right)^n}{n}\,.
\end{gather*}
Using \eqref{ineq:ratio_diff_pol_eps}, we find that 
\[
\left|\frac{\frac{d\pi_{\theta^{\varepsilon}}}{d\mu}(a|s)-\frac{d\pi_{\theta}}{d\mu}(a|s)}{\frac{d\pi_{\theta}}{d\mu}(a|s)}\right|\le \left(e^{\tilde{R}_\theta } + 1\right)\left(\varepsilon \|\theta'-\theta\|_{2}+\varepsilon^2 \exp\left(\|\theta'-\theta\|_{2}\right)\right)\frac{\tilde{\pi}_{\theta}(A|s)}{\tilde{\pi}_{\theta^{\varepsilon}}(A|s)}.
\]
Thus, by Lemma~\ref{lem:app:mdp:grad_of_log_policy}, we have
\begin{equation}
\underset{\varepsilon \in [0,1]}{\lim_{\varepsilon\rightarrow 0}}\frac{\log\frac{d\pi_{\theta^{\varepsilon}}}{d\mu}(a|s) -\log\frac{d\pi_{\theta}}{d\mu}(a|s) }{\varepsilon}= g(s,a) - \int g(s,a^\prime)\pi_\theta(da^\prime),\label{lem:func_deriv_log_pi}
\end{equation}
and that there exists a constant $C>0$ such that for all $\varepsilon<\varepsilon_0$, $s\in S$, and $\mu-a.e. \,a\in A$,
\[
\varepsilon^{-1}\left|\log\frac{d\pi_{\theta^{\varepsilon}}}{d\mu}(a|s) -\log\frac{d\pi_{\theta}}{d\mu}(a|s)\right| \le C\,.
\]
Therefore, owing to \eqref{eq:limit_policy_again} and \eqref{lem:func_deriv_log_pi}, we find
\[
\underset{\varepsilon \in [0,1]}{\lim_{\varepsilon\rightarrow 0}}I_4^{\varepsilon}(s)=+\tau\int\left( g(s,a)-\int g(s,a^\prime)\pi_{\theta}(da'|s)\right)\pi_{\theta}(da|s)=0\,.
\]
We now turn our attention to  $I_3^{\varepsilon}$. Recalling \eqref{eq:intro:value_function_occ}, we have 
\[
V^{\pi_{\theta^{\varepsilon}}}_{\tau}(s)=\frac{1}{1-\gamma}\int_{S}\int_A\left(c(s',a)+\tau\log\frac{d\pi_{\theta^{\varepsilon}}}{d\mu}(a|s')\right)\pi_{\theta^{\varepsilon}}(da|s')d^{\pi_{\theta^{\varepsilon}}}(ds'|s)\,.
\]

It follows from Corollary \ref{cor:app:mdp:Lip_occ}, Eq.(\ref{eq:limit_policy}) (also Eq.\eqref{eq:limit_policy_again}), Lemma \ref{lem:bounds}, and the bound on cost function, that 
\[
\underset{\varepsilon \in [0,1]}{\lim_{\varepsilon\rightarrow 0}}V^{\pi_{\theta^{\varepsilon}}}_{\tau}(s)=V^{\pi_{\theta}}_{\tau}(s)\,.
\]
Thus, by Lemmas \ref{lem:bounds} and \ref{lem:app:mdp:grad_of_log_policy}, by dominated convergence theorem we obtain 
\[
\underset{\varepsilon \in [0,1]}{\lim_{\varepsilon\rightarrow 0}}I_3^{\varepsilon}(s)=\gamma\int_A \int_S V^{\pi_{\theta}}_{\tau}(s')P(ds'|s,a)  \nabla \log \frac{d \pi_\theta}{d \mu}(a|s)\pi_{\theta}(da|s)\,.
\]
Using Lemmas \ref{lem:bounds} and \ref{lem:app:mdp:grad_of_log_policy} and bound on cost function, by dominated convergence theorem we get
\begin{equation*}
\underset{\varepsilon \in [0,1]}{\lim_{\varepsilon\rightarrow 0}}(I_1^{\varepsilon}(s)+I_5^{\varepsilon}(s))=\int_A \left(c(s',a)+\tau\log\frac{d\pi_{\theta}}{d\mu}(a|s)\right)\nabla \log \frac{d \pi_\theta}{d \mu}(a|s)\pi_{\theta}(da|s)\,.
\end{equation*}
Putting it all together and using the definition of $Q^{\pi}_\tau$, we arrive at
\[
\underset{\varepsilon \in [0,1]}{\lim_{\varepsilon\rightarrow 0}}\sum_{j=1}^5 I_j^{\varepsilon}(s)=\int_A \left(Q_{\tau}^{\pi_{\theta}}(s,a)+\tau\log\frac{d\pi_{\theta}}{d\mu}(a|s)\right)\nabla \log \frac{d \pi_\theta}{d \mu}\pi_{\theta}(da|s).
\]
Since $(I_j)_{1\le j\le 5}$ are bounded uniformly in $\varepsilon$, we may apply the bounded convergence theorem to pass to the limit in \eqref{eq:Vtau_diff_quot} to obtain \eqref{eq:wts_func_deriv_V}. 
\end{proof}

\subsubsection{Proof of Proposition \ref{prp:app:mdp:local_lipchitiz_conitinuity_mdp}}
\label{sec:proof_local_lip}

\begin{proof}
\begin{align}
&  \quad \ \nabla V_\tau^{\pi_\theta}(\rho) - \nabla V_\tau^{\pi_{\theta^\prime}}(\rho) \nonumber\\
& = \frac{1}{1-\gamma}\int_{S} \int_{A} \left(Q^{\pi_\theta}_\tau(s,a) + \tau \log \frac{d \pi_{\theta}}{d \mu}(a|s)\right)\nabla \log \frac{d\pi_\theta}{d \mu}(a| s)\pi_{\theta}(da|s)d_\rho^{\pi_{\theta}}(ds) \nonumber \\
& \quad - \frac{1}{1-\gamma}\int_{S} \int_{A} \left(Q^{\pi_{\theta^\prime}}_\tau(s,a) + \tau \log \frac{d \pi_{\theta^\prime}}{d \mu}(a|s)\right) \nabla \log \frac{d\pi_{\theta^\prime}}{d \mu}(a| s)\pi_{\theta^\prime}(da|s)d_\rho^{\pi_{\theta^\prime}}(ds) \nonumber\\
& = \frac{1}{1-\gamma}\int_{S} \int_{A} \left(Q^{\pi_\theta}_\tau(s,a) + \tau \log \frac{d \pi_{\theta}}{d \mu}(a|s) - Q^{\pi_{\theta^\prime}}_\tau(s,a) - \tau \log \frac{d \pi_{\theta^\prime}}{d \mu}(a|s)\right)\nonumber\\
& \qquad \qquad \qquad \ \cdot \nabla \log \frac{d\pi_\theta}{d \mu}(a| s)\pi_{\theta}(da|s)d_\rho^{\pi_{\theta}}(ds) \nonumber \\
& \quad +  \frac{1}{1-\gamma}\int_{S} \int_{A} \left(Q^{\pi_{\theta^\prime}}_\tau(s,a) + \tau \log \frac{d \pi_{\theta^\prime}}{d \mu}(a|s)\right)\nabla \log \frac{d\pi_{\theta}}{d \mu}(a| s)\pi_{\theta}(da|s)d_\rho^{\pi_{\theta}}(ds) \nonumber\\
& \quad - \frac{1}{1-\gamma}\int_{S} \int_{A} \left(Q^{\pi_{\theta^\prime}}_\tau(s,a) + \tau \log \frac{d \pi_{\theta^\prime}}{d \mu}(a|s)\right)\nabla \log \frac{d\pi_{\theta^\prime}}{d \mu}(a| s)\pi_{\theta^\prime}(da|s)d_\rho^{\pi_{\theta^\prime}}(ds) \nonumber\\
& = \frac{1}{1-\gamma}\int_{S} \int_{A} \left(Q^{\pi_\theta}_\tau(s,a) + \tau \log \frac{d \pi_{\theta}}{d \mu}(a|s) - Q^{\pi_{\theta^\prime}}_\tau(s,a) - \tau \log \frac{d \pi_{\theta^\prime}}{d \mu}(a|s)\right) \nonumber\\
& \qquad \qquad \qquad \ \cdot \nabla \log \frac{d\pi_\theta}{d \mu}(a| s)\pi_{\theta}(da|s)d_\rho^{\pi_{\theta}}(ds) \nonumber \\
& \quad + \frac{1}{1-\gamma}\int_{S} \int_{A} \left(Q^{\pi_{\theta^\prime}}_\tau(s,a) + \tau \log \frac{d \pi_{\theta^\prime}}{d \mu}(a|s)\right) \nonumber \\
& \qquad \qquad \qquad \ \cdot\left(\nabla \log \frac{d\pi_{\theta}}{d \mu}(a| s) - \nabla \log \frac{d\pi_{\theta^\prime}}{d \mu}(a| s)\right)\pi_{\theta}(da|s)d_\rho^{\pi_{\theta}}(ds) \nonumber\\
& \quad + \frac{1}{1-\gamma}\int_{S} \int_{A} \left(Q^{\pi_{\theta^\prime}}_\tau(s,a) + \tau \log \frac{d \pi_{\theta^\prime}}{d \mu}(a|s)\right)  \nabla \log \frac{d\pi_{\theta^\prime}}{d \mu}(a| s)\pi_{\theta}(da|s)d_\rho^{\pi_{\theta}}(ds) \nonumber\\
& \quad - \frac{1}{1-\gamma}\int_{S} \int_{A} \left(Q^{\pi_{\theta^\prime}}_\tau(s,a) + \tau \log \frac{d \pi_{\theta^\prime}}{d \mu}(a|s)\right)  \nabla \log \frac{d\pi_{\theta^\prime}}{d \mu}(a| s)\pi_{\theta^\prime}(da|s)d_\rho^{\pi_{\theta^\prime}}(ds) \nonumber\\
& = \frac{1}{1-\gamma}\int_{S} \int_{A} \left(Q^{\pi_\theta}_\tau(s,a) + \tau \log \frac{d \pi_{\theta}}{d \mu}(a|s) - Q^{\pi_{\theta^\prime}}_\tau(s,a) - \tau \log \frac{d \pi_{\theta^\prime}}{d \mu}(a|s)\right) \nonumber \\
& \qquad \qquad \qquad \ \cdot\nabla \log \frac{d\pi_\theta}{d \mu}(a| s)\pi_{\theta}(da|s)d_\rho^{\pi_{\theta}}(ds) \nonumber \\
& \quad + \frac{1}{1-\gamma}\int_{S} \int_{A} \left(Q^{\pi_{\theta^\prime}}_\tau(s,a) + \tau \log \frac{d \pi_{\theta^\prime}}{d \mu}(a|s)\right) \nonumber \\
& \qquad \qquad \qquad \ \cdot\left(\nabla \log \frac{d\pi_{\theta}}{d \mu}(a| s) - \nabla \log \frac{d\pi_{\theta^\prime}}{d \mu}(a| s)\right)\pi_{\theta}(da|s)d_\rho^{\pi_{\theta}}(ds) \nonumber\\
& \quad + \frac{1}{1-\gamma}\int_{S} \int_{A} \left(Q^{\pi_{\theta^\prime}}_\tau(s,a) + \tau \log \frac{d \pi_{\theta^\prime}}{d \mu}(a|s)\right)  \nonumber \\
& \qquad \qquad \qquad \ \cdot \nabla \log \frac{d\pi_{\theta^\prime}}{d \mu}(a| s)\left(\pi_{\theta}(da|s) -\pi_{\theta^\prime}(da|s)\right)d_\rho^{\pi_{\theta}}(ds) \nonumber\\
& \quad + \frac{1}{1-\gamma}\int_{S} \int_{A} \left(Q^{\pi_{\theta^\prime}}_\tau(s,a) + \tau \log \frac{d \pi_{\theta^\prime}}{d \mu}(a|s)\right)  \nonumber \\
& \qquad \qquad \qquad \ \cdot
\nabla \log \frac{d\pi_{\theta^\prime}}{d \mu}(a| s)\pi_{\theta^\prime}(da|s)\left(d_\rho^{\pi_{\theta}}(ds)- d_\rho^{\pi_{\theta^\prime}}(ds)\right), \label{eq: difference of the derivative}
\end{align}
where 
\begin{align}
& \frac{1}{1-\gamma}\int_{S} \int_{A} \left(Q^{\pi_\theta}_\tau(s,a) + \tau \log \frac{d \pi_{\theta}}{d \mu}(a|s) - Q^{\pi_{\theta^\prime}}_\tau(s,a) - \tau \log \frac{d \pi_{\theta^\prime}}{d \mu}(a|s)\right) \nonumber \\
& \qquad \qquad \ \ \ \cdot\nabla \log \frac{d\pi_\theta}{d \mu}(a| s)\pi_{\theta}(da|s)d_\rho^{\pi_{\theta}}(ds)  \nonumber \\
& =  \frac{1}{1-\gamma}\int_{S} \int_{A} \bigg(\gamma \int_{S} \left(V^{\pi_\theta}_\tau (s^\prime) - V^{\pi_{\theta^\prime}}_\tau (s^\prime)\right)P(d s^\prime | s,a) \nonumber \\
& \quad \quad \qquad \qquad \ + \tau \langle g(s,a), \theta - \theta^\prime \rangle + \tau\left(Z_{\pi_{\theta^\prime}}(s) - Z_{\pi_\theta}(s)\right)\bigg) \nonumber\\
& \qquad \qquad \qquad \ \cdot\nabla \log \frac{d\pi_\theta}{d \mu}(a| s)\pi_{\theta}(da|s)d_\rho^{\pi_{\theta}}(ds) \nonumber\\
& =  \frac{1}{1-\gamma}\int_{S} \int_{A} \left(\gamma \int_{S}  \int_0^1 \left\langle \nabla V_\tau^{\pi_{\theta^\epsilon}}(\rho), \theta - \theta^\prime\right\rangle d \epsilon P(d s^\prime | s,a) + \tau \langle g(s,a), \theta - \theta^\prime \rangle \right) \nonumber\\
& \qquad \qquad \qquad \ \cdot\nabla \log \frac{d\pi_\theta}{d \mu}(a| s)\pi_{\theta}(da|s)d_\rho^{\pi_{\theta}}(ds) \label{eq: difference of the derivative 01} 
\end{align}

Using Lemma \ref{lem:app:mdp:bounded_grad_of_log_policy}, Lemma \ref{lem:bounds}, Lemma \ref{lem: smooth of log policy}, Corollary \ref{cor:app:mdp:Lip_occ}, putting Eq.(\ref{eq: difference of the derivative 01}) back to Eq.(\ref{eq: difference of the derivative}), we have 
\begin{align}
& \left\|\nabla V_\tau^{\pi_\theta}(\rho) - \nabla V_\tau^{\pi_{\theta^\prime}}(\rho)\right\|_2\\
& \leq \left( \frac{1}{1-\gamma}\left(\frac{ \gamma (5 + \tau R)  }{1-\gamma} + 6\right)\left(\frac{1+\gamma\tau R}{1-\gamma} + \tau R\right) + \frac{2\tau }{1-\gamma}\right) \Bigg\|\theta - \theta^\prime\Bigg\|_2.
\end{align}
\end{proof}

\section{Conclusion}

We prove linear convergence of policy gradient for entropy regularized MDPs with log-linear policies on general state and action spaces under $Q^\pi_\tau$-realizability and when suitable basis functions are employed. 
This complements existing results for softmax policy gradient methods in the tabular setting~\cite{mei2020global}.
To obtain our results we have established a non-uniform P{\L} inequality for general state and action spaces and carried out novel Lyapunov function-based analysis allowing control of the non-uniform term.   


\subsection*{Acknowledgements}
The second author was partially supported by a grant from the Simons Foundation.
The second author would also like to thank the Isaac Newton Institute for Mathematical Sciences, Cambridge, for support and hospitality during the programme Bridging Stochastic Control And Reinforcement Learning: Theories and Applications, where work on this paper was partially undertaken. This work was supported by EPSRC grant EP/V521929/1.
The second and third authors acknowledge funding from the UKRI Prosperity Partnerships grant APP43592: AI$^2$ - Assurance and Insurance for Artificial Intelligence, which supported this work.



\bibliographystyle{plainnat}
\bibliography{references} 



\end{document}
